%==============================================================================
% BookNest --- Complete Technical Guide & Demo Preparation Handbook
% Team Neuron
%
% Build with:  tectonic docs/booknest_demo_guide.tex
%          or: pdflatex docs/booknest_demo_guide.tex   (run twice for TOC)
%==============================================================================
\documentclass[11pt,a4paper]{article}

\usepackage[margin=2.3cm,bottom=2.6cm]{geometry}
\usepackage[T1]{fontenc}
\usepackage{lmodern}
\usepackage{xcolor}
\usepackage{titlesec}
\usepackage{enumitem}
\usepackage{fancyhdr}
\usepackage{graphicx}
\usepackage{booktabs}
\usepackage{longtable}
\usepackage{array}
\usepackage{listings}
\usepackage{parskip}
\usepackage{mdframed}
\usepackage[hidelinks]{hyperref}

%------------------------------------------------------------------ colours
\definecolor{primary}{RGB}{15, 90, 70}
\definecolor{gold}{RGB}{201, 162, 77}
\definecolor{codebg}{RGB}{247, 247, 245}
\definecolor{codeline}{RGB}{215, 215, 210}
\definecolor{notebg}{RGB}{239, 246, 255}
\definecolor{noteborder}{RGB}{59, 130, 246}
\definecolor{warnbg}{RGB}{255, 247, 237}
\definecolor{warnborder}{RGB}{234, 88, 12}
\definecolor{saybg}{RGB}{240, 253, 244}
\definecolor{sayborder}{RGB}{22, 163, 74}
\definecolor{kw}{RGB}{155, 35, 147}
\definecolor{cmt}{RGB}{110, 120, 130}
\definecolor{str}{RGB}{196, 26, 22}

\hypersetup{colorlinks=true, linkcolor=primary, urlcolor=primary,
            pdftitle={BookNest Technical Guide}, pdfauthor={Team Neuron}}

%------------------------------------------------------------------ headings
\titleformat{\section}{\Large\bfseries\color{primary}}{\thesection.}{0.6em}{}[{\titlerule[0.8pt]}]
\titleformat{\subsection}{\large\bfseries}{\thesubsection}{0.6em}{}
\titleformat{\subsubsection}{\normalsize\bfseries\color{primary}}{}{0em}{}
\setlength{\parskip}{6pt}

%------------------------------------------------------------------ listings
\lstdefinestyle{code}{
  backgroundcolor=\color{codebg},
  basicstyle=\ttfamily\scriptsize,
  keywordstyle=\color{kw}\bfseries,
  commentstyle=\color{cmt}\itshape,
  stringstyle=\color{str},
  breaklines=true,
  breakatwhitespace=false,
  columns=fullflexible,
  keepspaces=true,
  showstringspaces=false,
  frame=single,
  rulecolor=\color{codeline},
  framesep=5pt,
  xleftmargin=4pt,
  xrightmargin=2pt,
  aboveskip=8pt,
  belowskip=8pt,
  captionpos=b,
  numbers=none,
  tabsize=2,
}
\lstset{style=code}

\lstdefinelanguage{TS}{
  morekeywords={const,let,var,function,async,await,return,if,else,for,while,
    import,export,from,type,interface,new,class,extends,implements,try,catch,
    finally,throw,default,of,in,null,undefined,true,false,this,as,void,Promise},
  sensitive=true,
  morecomment=[l]{//},
  morecomment=[s]{/*}{*/},
  morestring=[b]",
  morestring=[b]',
  morestring=[b]`,
}

%------------------------------------------------------------------ boxes
\newmdenv[backgroundcolor=notebg, linecolor=noteborder, linewidth=2.5pt,
  leftline=true, rightline=false, topline=false, bottomline=false,
  innerleftmargin=9pt, innerrightmargin=9pt, innertopmargin=7pt,
  innerbottommargin=7pt, skipabove=8pt, skipbelow=8pt]{notebox}

\newmdenv[backgroundcolor=warnbg, linecolor=warnborder, linewidth=2.5pt,
  leftline=true, rightline=false, topline=false, bottomline=false,
  innerleftmargin=9pt, innerrightmargin=9pt, innertopmargin=7pt,
  innerbottommargin=7pt, skipabove=8pt, skipbelow=8pt]{warnbox}

\newmdenv[backgroundcolor=saybg, linecolor=sayborder, linewidth=2.5pt,
  leftline=true, rightline=false, topline=false, bottomline=false,
  innerleftmargin=9pt, innerrightmargin=9pt, innertopmargin=7pt,
  innerbottommargin=7pt, skipabove=8pt, skipbelow=8pt]{saybox}

% "Say this in the demo" helper
\newcommand{\saythis}[1]{\begin{saybox}\textbf{\small\color{primary}SAY THIS:} \textit{#1}\end{saybox}}

% FAQ helpers
\newcommand{\Q}[1]{\vspace{7pt}\noindent\textbf{\color{primary}Q.}~\textbf{#1}\par\vspace{1pt}}
\newcommand{\A}[1]{\noindent #1\par}

\newcommand{\file}[1]{\texttt{\small #1}}

%------------------------------------------------------------------ header
\pagestyle{fancy}
\fancyhf{}
\rhead{\small\color{primary}BookNest --- Technical Guide}
\lhead{\small Team Neuron}
\cfoot{\small\thepage}
\renewcommand{\headrulewidth}{0.4pt}

\begin{document}

%==============================================================================
\begin{titlepage}
\centering
\vspace*{2.2cm}
{\Huge\bfseries\color{primary} BookNest\par}
\vspace{0.4cm}
{\large Hotel Management \& Booking System\par}
\vspace{1.4cm}
\rule{0.75\textwidth}{1.2pt}\par
\vspace{0.7cm}
{\LARGE\bfseries Complete Technical Guide\par}
\vspace{0.35cm}
{\large \& Demo Preparation Handbook\par}
\vspace{0.7cm}
\rule{0.75\textwidth}{1.2pt}\par
\vspace{1.4cm}

\begin{minipage}{0.82\textwidth}
\centering\itshape
Every system in this project, explained from first principles ---
what it does, how it is wired, why we chose it, and the exact
answers to the questions you will be asked.
\end{minipage}

\vspace{2.4cm}
{\large\bfseries Team Neuron\par}
\vspace{0.3cm}
{\small \href{https://hms-and-booking.vercel.app}{hms-and-booking.vercel.app}\par}
\vspace{1.6cm}
{\small \today\par}
\end{titlepage}

%==============================================================================
\tableofcontents
\newpage

%==============================================================================
\section{How To Use This Document}

This document has three kinds of content, and you should read them differently.

\begin{enumerate}[leftmargin=*]
  \item \textbf{Explanations.} Plain-English descriptions of every system. Read
        these to \emph{understand}. If a sentence uses a word you cannot define,
        the Glossary (Section~\ref{sec:glossary}) defines it.
  \item \textbf{Code.} The actual code from the repository, with the file path
        above it. You should be able to point at any of these and explain what
        each line does.
  \item \textbf{Green ``SAY THIS'' boxes.} A one-or-two-sentence answer, written
        the way you would speak it out loud. If you are asked a question and your
        mind goes blank, say the green box.
\end{enumerate}

\begin{notebox}
\textbf{The single most important piece of advice.} When someone asks
``how does X work?'', do not start with the code. Start with the \emph{problem}
X solves, then say what you built, then offer the detail. For example:
\emph{``The problem is two people can book the last room at the same moment.
So booking is not done in JavaScript --- it is a Postgres function that takes a
row lock on each night's inventory before it writes. Want me to walk through it?''}
That sentence alone tells an examiner you understand your own system.
\end{notebox}

\newpage
%==============================================================================
\section{The Big Picture}

\subsection{What BookNest is}

BookNest is a hotel booking platform. It is not a toy: it takes real payments
through a real payment gateway, sends real email, and enforces real security
rules at the database level. It has four kinds of users:

\begin{longtable}{@{}p{2.4cm}p{12.4cm}@{}}
\toprule
\textbf{Role} & \textbf{What they can do} \\
\midrule
\endhead
\textbf{Guest} & Search hotels, view listings, book rooms, pay online, check in
with a QR code, review their stay, message the hotel. \\
\addlinespace
\textbf{Manager} & Create hotel listings through a 9-step wizard, run a front
desk (check guests in/out, take walk-in bookings), view earnings, invite staff,
reply to guest messages. \\
\addlinespace
\textbf{Staff} & A real login account invited by a manager. Gets a subset of the
manager's powers on \emph{one specific hotel}, controlled by a permissions array. \\
\addlinespace
\textbf{Admin} & Approve or reject manager verification requests and hotel
listings, suspend users, change roles, view platform statistics. Every admin
action is written to an audit log. \\
\bottomrule
\end{longtable}

\subsection{The technology, in one line each}

\begin{longtable}{@{}p{3.5cm}p{11.3cm}@{}}
\toprule
\textbf{Layer} & \textbf{Choice and one-line justification} \\
\midrule
\endhead
Framework & \textbf{Next.js 16} (App Router). One codebase serves both the
website and the API, so we never maintain two projects. \\
\addlinespace
UI library & \textbf{React 19}. Industry standard; Next.js is built on it. \\
\addlinespace
Language & \textbf{TypeScript}. Catches type mistakes before the code runs. \\
\addlinespace
Database & \textbf{Supabase} = hosted \textbf{PostgreSQL} + Auth + Storage +
Realtime. One service instead of four. \\
\addlinespace
Security model & \textbf{Row Level Security} (Postgres). The database itself
refuses to return rows you are not allowed to see. \\
\addlinespace
Payments & \textbf{Razorpay}. India-first, supports UPI, cards, netbanking.
Has a proper test mode. \\
\addlinespace
Realtime chat & \textbf{Supabase Realtime} (WebSockets over Postgres
replication). No separate socket server to run. \\
\addlinespace
Maps & \textbf{MapLibre GL} + \textbf{OpenFreeMap}, with \textbf{Leaflet}
as a fallback. Zero API keys, zero cost, no billing surprises. \\
\addlinespace
Email & \textbf{Brevo} HTTP API. Works from serverless functions (SMTP often
does not) and has a free tier without needing a custom domain. \\
\addlinespace
Styling & \textbf{Tailwind CSS v4}. Styles live next to the markup. \\
\addlinespace
Hosting & \textbf{Vercel}. Made by the Next.js team; push to GitHub and it deploys. \\
\bottomrule
\end{longtable}

\subsection{The life of a request}

This diagram is worth memorising. Almost every question about ``how does the app
work'' is answered by pointing at one box in it.

\begin{lstlisting}[basicstyle=\ttfamily\scriptsize]
                            BROWSER
                               |
        +----------------------+----------------------+
        |                      |                      |
    (1) Page load        (2) Data fetch        (3) Live updates
        |                      |                      |
        v                      v                      v
  +-----------+        +---------------+     +------------------+
  | proxy.ts  |        | Route Handler |     | Supabase Realtime|
  | (gate)    |        | /api/...      |     | WebSocket        |
  +-----+-----+        +-------+-------+     +---------+--------+
        |                      |                       |
        v                      v                       |
  +-----------------------------------+                |
  |  Next.js Server Component (RSC)   |                |
  |  - runs ONLY on the server        |                |
  |  - can read the DB directly       |                |
  +-----------------+-----------------+                |
                    |                                  |
                    v                                  |
             +-------------------------------+         |
             |      SUPABASE (PostgreSQL)    |<--------+
             |  RLS policies + RPC functions |
             |  Auth  |  Storage  | Realtime |
             +-------------------------------+
                    ^                    ^
                    |                    |
             Razorpay webhook      Brevo email API
             (POST /api/webhooks)  (outbound HTTPS)
\end{lstlisting}

\saythis{The browser talks to Next.js. Next.js talks to Postgres. But even if
somebody bypassed Next.js entirely and hit the database directly with our public
key, Row Level Security would still stop them from reading someone else's
bookings. Security is enforced at the bottom of the stack, not the top.}

\newpage
%==============================================================================
\section{Next.js: Structure and How Files Are Linked}

\subsection{The one rule that explains the whole folder}

In the Next.js \textbf{App Router}, \emph{a folder is a URL, and a file named
\file{page.tsx} is what the user sees at that URL.}

\begin{lstlisting}
app/page.tsx                -> /
app/hotels/page.tsx         -> /hotels
app/hotels/[id]/page.tsx    -> /hotels/<any-id>       ([id] = dynamic segment)
app/manager/staff/page.tsx  -> /manager/staff
app/api/bookings/route.ts   -> /api/bookings          (an API endpoint, not a page)
\end{lstlisting}

There is no router configuration file anywhere in this project. There is no
\file{routes.js}. The folder structure \emph{is} the routing table. This is why
the \file{app/} directory looks the way it does.

\subsection{The special filenames}

Inside any folder, certain filenames have reserved meaning. These are the ones
we use:

\begin{longtable}{@{}p{2.9cm}p{11.9cm}@{}}
\toprule
\textbf{Filename} & \textbf{What Next.js does with it} \\
\midrule
\endhead
\file{page.tsx} & The page itself. Must be present for the URL to exist. \\
\addlinespace
\file{layout.tsx} & A wrapper that persists across navigations inside that
folder. \file{app/layout.tsx} is the root layout (html/body, fonts, navbar).
\file{app/manager/layout.tsx} wraps every manager page in the sidebar shell. \\
\addlinespace
\file{loading.tsx} & Shown automatically while \file{page.tsx} is fetching data.
Next.js wraps the page in a React \texttt{<Suspense>} boundary for you. This is
why we have \file{app/hotels/loading.tsx} showing skeleton cards. \\
\addlinespace
\file{template.tsx} & Like a layout, but re-created on every navigation. We use
\file{app/manager/template.tsx} and \file{app/admin/template.tsx} for page
transition animations, which need to re-mount to re-run. \\
\addlinespace
\file{route.ts} & Turns the folder into an \textbf{API endpoint} instead of a
page. Exports functions named \texttt{GET}, \texttt{POST}, etc. \\
\addlinespace
\file{actions.ts} & Not a Next.js reserved name --- \emph{our} convention. Holds
Server Actions (see below). Marked with \texttt{"use server"} at the top. \\
\addlinespace
\file{proxy.ts} & Runs \emph{before} every matching request. In Next.js 15 and
earlier this file was called \file{middleware.ts}; Next.js 16 renamed the
convention to \file{proxy.ts}. Ours guards the protected routes. \\
\bottomrule
\end{longtable}

\begin{notebox}
\textbf{Why is our middleware called \file{proxy.ts}?} Because we are on Next.js
16, which renamed the convention. The function inside is still just a function
that receives a request and returns a response. If an examiner says ``I thought
middleware was called middleware.ts'', they are right --- for Next.js 15. This is
a good thing to know; it shows you are on the current version and understand why.
\end{notebox}

\subsection{Server Components vs Client Components --- the mental model}

This is the concept most people get wrong, so learn it properly.

\textbf{By default, every component in \file{app/} is a Server Component.} It
runs on the server, once, and only its resulting HTML is sent to the browser.
Its JavaScript is \emph{never downloaded by the user}. It can \texttt{await}
directly. It can read the database. It can read secrets.

\textbf{A file that starts with \texttt{"use client"} is a Client Component.}
Its JavaScript \emph{is} shipped to the browser. It can use \texttt{useState},
\texttt{useEffect}, \texttt{onClick}, and browser APIs like the camera. It
cannot read secrets, because the user can read its source.

\begin{lstlisting}[language=TS,caption={app/hotels/[id]/page.tsx --- a Server Component. Note: no "use client", and it awaits a database read directly.}]
export const revalidate = 60;   // ISR: regenerate this page at most once a minute

export default async function HotelDetailPage({ params }) {
  const { id } = await params;

  // This runs on the SERVER. The user never sees this query.
  let detail = await getPublicHotel(id);
  if (!detail) notFound();

  return <div> ... </div>;
}
\end{lstlisting}

The rule of thumb we followed across the whole codebase:

\begin{quote}
\emph{Fetch data in a Server Component. Pass it down as props into a small
Client Component that handles the interactivity.}
\end{quote}

You can see this exact pattern everywhere:

\begin{itemize}[leftmargin=*,noitemsep]
  \item \file{app/dashboard/page.tsx} (server, fetches) renders
        \file{app/dashboard/DashboardClient.tsx} (client, interactive).
  \item \file{app/hotels/page.tsx} renders \file{app/hotels/HotelsClient.tsx}.
  \item \file{app/manager/staff/page.tsx} renders \file{StaffClient.tsx}.
  \item \file{app/messages/page.tsx} renders \file{GuestMessagesClient.tsx}.
\end{itemize}

\saythis{Server Components render on the server and ship zero JavaScript, so the
page is fast and the database query never reaches the browser. We only opt into
a Client Component where we genuinely need interactivity --- a map, a date picker,
the chat box. That is why the naming convention across the repo is
\texttt{page.tsx} for the server half and \texttt{SomethingClient.tsx} for the
interactive half.}

\subsection{The four ways data moves in this app}

We do not use one data-fetching strategy. We deliberately use four, each where
it fits. Being able to say \emph{why} each one is used where it is used is the
single most impressive thing you can do in this demo.

\subsubsection{1. Server Component reads the database directly}

Used for: public pages that must be fast and SEO-friendly (home page, hotel
list, hotel detail page).

\begin{lstlisting}[language=TS]
// app/hotels/[id]/page.tsx  -- runs on server, no API route involved
const detail = await getPublicHotel(id);   // lib/hotels.ts -> Supabase
\end{lstlisting}

No API endpoint exists for this. The page talks to Postgres and returns HTML.
One network round trip for the user instead of two.

\subsubsection{2. Route Handler + SWR (client-side fetching)}

Used for: dashboards, where data changes and the user expects it to refresh.

\begin{lstlisting}[language=TS,caption={lib/swr.ts --- one shared fetcher so preloading and fetching use the same cache key}]
export const fetcher = (url: string) =>
  fetch(url).then((r) => {
    if (!r.ok) throw new Error(`Request to ${url} failed with status ${r.status}`);
    return r.json();
  });

// Maps a manager nav route to the API endpoint that page reads, for preloading.
export function managerApiForRoute(href: string): string | null {
  if (href.startsWith("/manager/hotels")) return "/api/manager/hotels";
  if (href.startsWith("/manager/manage")) return "/api/manager/manage";
  if (href.startsWith("/manager/staff"))  return "/api/manager/staff";
  return null;
}
\end{lstlisting}

\textbf{SWR} stands for \emph{stale-while-revalidate}. When you open a page, SWR
instantly shows the last data it had (stale), then quietly re-fetches in the
background and swaps in the fresh data. The dashboard feels instant.

The \texttt{managerApiForRoute} function powers a nice trick: when you
\emph{hover} over a sidebar link, we call SWR's \texttt{preload()} for the
endpoint that page will need. By the time you click, the data has already
arrived.

\subsubsection{3. Server Actions}

Used for: writes triggered by a form or a button, where we want to avoid writing
an API endpoint at all.

A Server Action is a function marked \texttt{"use server"} that you can
\emph{import and call directly from a Client Component}. Next.js secretly turns
the call into an HTTP POST. You never write the fetch, the URL, or the JSON.

\begin{lstlisting}[language=TS,caption={app/messages/actions.ts}]
"use server";

export async function markRead(conversationId: string) {
  const supabase = await createClient();
  const { data: { user } } = await supabase.auth.getUser();
  if (!user) return;

  const { data: conversation } = await supabase
    .from("conversations").select("guest_id").eq("id", conversationId).single();
  if (!conversation) return;

  const isGuest = conversation.guest_id === user.id;
  await supabase
    .from("conversations")
    .update(isGuest ? { guest_unread: 0 } : { host_unread: 0 })
    .eq("id", conversationId);

  revalidatePath(`/messages`);          // tell Next.js those pages are stale
  revalidatePath(`/manager/messages`);
}
\end{lstlisting}

In the chat hook we simply write \texttt{await apiMarkRead(id)}. It looks like a
local function call. It is actually a secure server round trip.

\begin{warnbox}
\textbf{The trap.} A Server Action is a public HTTP endpoint whether you like it
or not. Anyone can call it with any arguments. That is why \emph{every} action in
this codebase re-checks \texttt{auth.getUser()} and relies on RLS --- we never
assume the caller is who the UI says they are.
\end{warnbox}

\subsubsection{4. Realtime WebSocket subscription}

Used for: chat only. The browser opens a persistent connection to Supabase and
Postgres pushes new rows to it. Covered fully in Section~\ref{sec:messaging}.

\subsection{Caching: how the public pages stay fast}

Three layers, all visible in the code.

\textbf{(a) \texttt{unstable\_cache} with a tag.} The approved-hotels query is
expensive (it joins rooms, reviews and photos). We cache its \emph{result} in
Next.js's data cache, under the tag \texttt{"hotels"}.

\begin{lstlisting}[language=TS,caption={lib/hotels.ts}]
export const HOTELS_CACHE_TAG = "hotels";

export const getApprovedHotelsCached = unstable_cache(
  async (): Promise<ExploreHotel[]> => {
    const supabase = createPublicClient();
    const { data, error } = await supabase
      .from("hotels")
      .select(`id, name, location, ..., rooms (price), reviews (rating),
               hotel_photos (url, category, sort_order)`)
      .eq("status", "approved")
      .order("created_at", { ascending: false });
    if (error) throw error;
    return data ?? [];
  },
  ["approved-hotels-explore"],                 // cache key
  { tags: [HOTELS_CACHE_TAG], revalidate: 3600 }, // tag + 1 hour fallback
);
\end{lstlisting}

\textbf{(b) Tag invalidation.} When an admin approves a hotel, or a manager
edits one, we call \texttt{revalidateTag("hotels")} (see
\file{app/actions/hotels.ts} and \file{app/api/clear-cache/route.ts}). Every
cached query carrying that tag is thrown away instantly. So the cache is one
hour \emph{or until something actually changes}, whichever comes first.

\textbf{(c) ISR on the page itself.} \file{app/hotels/[id]/page.tsx} exports
\texttt{revalidate = 60}. The rendered HTML is stored on Vercel's CDN and
regenerated at most once a minute.

\saythis{The hotel page is cached at two levels: the database query result is
cached under a tag, and the rendered HTML is cached on the CDN for 60 seconds.
When an admin approves a listing we call \texttt{revalidateTag}, which
invalidates the query cache immediately rather than waiting for the timer. So
the site is fast \emph{and} correct.}

\subsection{\file{proxy.ts}: the gate in front of everything}

\begin{lstlisting}[language=TS,caption={proxy.ts}]
export async function proxy(request: NextRequest) {
  return updateSession(request);
}

export const config = {
  // Only run on the protected route groups. Public pages skip middleware
  // entirely -- no auth round trip. API routes authenticate themselves.
  matcher: ["/dashboard/:path*", "/manager/:path*", "/admin/:path*"],
};
\end{lstlisting}

The \texttt{matcher} is a deliberate performance decision. Middleware runs
before \emph{every} matching request, so if we matched \texttt{/} and
\texttt{/hotels} too, every anonymous visitor to the home page would pay for an
auth check they do not need. We match only the three private areas.

The real logic lives in \file{lib/supabase/middleware.ts}:

\begin{lstlisting}[language=TS,caption={lib/supabase/middleware.ts --- the access-control cascade}]
const isGuestArea   = path.startsWith("/dashboard");
const isManagerArea = path.startsWith("/manager");
const isAdminArea   = path.startsWith("/admin");
const isProtected   = isGuestArea || isManagerArea || isAdminArea;

// getClaims() verifies the JWT locally -- no round trip to the Auth server.
const { data: claims } = await supabase.auth.getClaims();
const userId = claims?.claims?.sub ?? null;

if (!userId) {
  if (isProtected) return redirect("/login?redirect=" + path);
  return response;                       // public page: cost nothing
}
if (!isProtected) return response;

// Role + (for managers) verification status, fetched IN PARALLEL.
const [{ data: profile }, { data: mv }] = await Promise.all([
  supabase.from("profiles").select("role, suspended").eq("id", userId).single(),
  isManagerArea
    ? supabase.from("manager_verifications").select("status")
        .eq("user_id", userId).order("created_at", { ascending: false })
        .limit(1).maybeSingle()
    : Promise.resolve({ data: null }),
]);

if (profile?.suspended)                        return redirectTo("/", "?suspended=1");
if (isAdminArea   && role !== "admin")         return redirectTo("/");
if (isGuestArea   && role !== "guest")         return redirectTo("/");
if (isManagerArea && role !== "manager" && role !== "staff") return redirectTo("/");

// Managers (not staff) are gated behind admin approval.
if (isManagerArea && role === "manager") {
  const onWaiting = path.startsWith("/manager/waiting");
  if (mv?.status !== "approved" && !onWaiting) return redirectTo("/manager/waiting");
  if (mv?.status === "approved" && onWaiting)  return redirectTo("/manager/hotels");
}
\end{lstlisting}

Three things to point out if asked:

\begin{enumerate}[leftmargin=*,noitemsep]
  \item \texttt{getClaims()} not \texttt{getUser()}. \texttt{getUser()} makes a
        network call to Supabase's auth server on \emph{every request}.
        \texttt{getClaims()} verifies the JWT signature locally using the
        project's public signing key. Same guarantee, no round trip.
  \item The two queries run inside \texttt{Promise.all}, not one after the other.
  \item Suspension is checked before role, so a suspended admin is still locked out.
\end{enumerate}

\begin{warnbox}
\textbf{Middleware is a convenience, not the security boundary.} If middleware
were the only protection, someone could hit \texttt{/api/admin/users} directly
and bypass it (\texttt{/api} is not in the matcher). That is fine, because every
admin API route independently calls \texttt{getAdminContext()}, \emph{and} RLS
would refuse the query anyway. Middleware exists to give users a nice redirect,
not to keep attackers out.
\end{warnbox}

\subsection{The complete file map}

\begin{lstlisting}
HMS-and-booking/
|
|-- proxy.ts .................. Auth/role gate. Runs before /dashboard, /manager, /admin
|-- next.config.ts ........... Image host allowlist
|-- app/
|   |-- layout.tsx ........... Root layout: fonts, AuthProvider, CurrencyProvider, Navbar
|   |-- page.tsx ............. Landing page
|   |-- globals.css .......... Tailwind + design tokens (brand colours)
|   |
|   |-- hotels/
|   |   |-- page.tsx ......... Server: fetches approved hotels (cached)
|   |   |-- HotelsClient.tsx . Client: filters, grid/map toggle
|   |   |-- ExploreMap.tsx ... Client: MapLibre pins
|   |   |-- loading.tsx ...... Skeleton shown while page.tsx fetches
|   |   `-- [id]/
|   |       |-- page.tsx ..... Server: getPublicHotel(), ISR revalidate=60
|   |       |-- components/ .. GalleryModal, BookingWidget, ReviewsSection, ...
|   |       `-- book/page.tsx  The booking flow
|   |
|   |-- bookings/[id]/
|   |   |-- page.tsx ......... Booking detail + generates the check-in QR (server-side)
|   |   |-- check-in/ ........ Guest-facing QR display
|   |   `-- review/ .......... Leave a review after checkout
|   |
|   |-- messages/ ............ Guest inbox      (+ actions.ts: server actions)
|   |-- dashboard/ ........... Guest bookings   (role=guest only)
|   |
|   |-- manager/
|   |   |-- layout.tsx ....... Sidebar shell (ManagerShell)
|   |   |-- create-hotel/ .... 9-step listing wizard (steps/Step1..Step9)
|   |   |-- manage/[hotelId]/  FrontDesk: QR scanner, check-in/out, walk-in bookings
|   |   |-- staff/ ........... Invite + permission management
|   |   `-- waiting/ ......... Shown until an admin approves the manager
|   |
|   |-- admin/ ............... Dashboard, verifications, listings, users
|   |
|   |-- auth/callback/route.ts  OAuth + email-confirmation landing; routes by role
|   |
|   `-- api/                    <-- every server endpoint
|       |-- bookings/ ........ GET list, POST create
|       |-- availability/ .... GET rooms free for a date range
|       |-- payments/razorpay/{order,verify,reconcile}
|       |-- webhooks/razorpay/  Razorpay calls THIS (server-authoritative)
|       |-- manager/{hotels,manage,staff}
|       |-- admin/{overview,users,listings,verifications}
|       `-- messages/{attachment,unread-count,...}
|
|-- lib/
|   |-- supabase/
|   |   |-- client.ts ........ Browser client   (anon key, cookies)
|   |   |-- server.ts ........ Server client    (anon key, request cookies)
|   |   |-- public.ts ........ Cookie-free client (safe inside unstable_cache)
|   |   |-- admin.ts ......... SERVICE ROLE - bypasses RLS. Webhooks only.
|   |   `-- middleware.ts .... The role cascade used by proxy.ts
|   |-- booking.ts ........... Price math, refund policy, date helpers
|   |-- razorpay.ts .......... Server SDK instance (holds the secret)
|   |-- razorpayCheckout.ts .. Browser: order -> checkout modal -> verify
|   |-- reconcile.ts ......... Safety net: ask Razorpay "was this actually paid?"
|   |-- email.ts ............. Brevo API + branded HTML layout
|   |-- emails/*.ts .......... One file per transactional email
|   |-- hotels.ts ............ Cached hotel queries + listing assembly
|   |-- compression.ts ....... Canvas-based client-side image compression
|   |-- auth.ts / authServer.ts
|   `-- types.ts ............. Every shared TypeScript type
|
|-- components/ .............. Shared UI (Navbar, AuthProvider, DatePicker, ...)
|   |-- messaging/ ........... ChatThread, Composer, useConversationMessages
|   |-- manager/ ............. ManagerShell, Panel
|   `-- admin/ ............... AdminShell
|
`-- supabase/ ................ 16 .sql migration files -- the real source of truth
\end{lstlisting}

\newpage
%==============================================================================
\section{Supabase: How the Database Is Implemented}

\subsection{What Supabase actually is}

Supabase is \textbf{a PostgreSQL database} with four services bolted on:

\begin{itemize}[leftmargin=*,noitemsep]
  \item \textbf{Auth} --- user accounts, passwords, Google OAuth, JWTs.
  \item \textbf{PostgREST} --- automatically turns every table into a REST API.
  \item \textbf{Storage} --- an S3-like file store, with permissions.
  \item \textbf{Realtime} --- streams Postgres changes over WebSockets.
\end{itemize}

\begin{notebox}
\textbf{The idea that makes Supabase different.} Because the database is exposed
directly to the browser through PostgREST, the browser holds a real database
API key (the ``anon key''). That sounds terrifying until you understand
\textbf{Row Level Security}: the database is told, per table, exactly which rows
each user may see. The anon key gets you \emph{to} the database; RLS decides what
you get \emph{out of} it. The key is public by design.
\end{notebox}

\subsection{Why there are four Supabase clients}
\label{sec:fourclients}

This is a question you will almost certainly be asked, because it looks like
duplication. It is not. Each client exists because of a different constraint.

\begin{longtable}{@{}p{2.8cm}p{2.1cm}p{9.6cm}@{}}
\toprule
\textbf{File} & \textbf{Key} & \textbf{Why it exists} \\
\midrule
\endhead
\file{client.ts} & anon & Runs in the browser. Reads and writes the auth cookie
so the session survives a refresh. \\
\addlinespace
\file{server.ts} & anon & Runs in Server Components, Route Handlers and Server
Actions. Reads the \emph{current request's} cookies via \texttt{cookies()}, so
queries run as the logged-in user and RLS applies to them. \\
\addlinespace
\file{public.ts} & anon & Stateless, \emph{cookie-free}. Needed because
\texttt{unstable\_cache} forbids touching request-bound data like cookies --- a
cached value must not depend on who asked for it. RLS still applies (anon can
only read approved hotels). \\
\addlinespace
\file{admin.ts} & \textbf{service role} & \textbf{Bypasses RLS entirely.} Used in
exactly one situation: the Razorpay webhook, where there is no logged-in user
because Razorpay's server is the caller. The service-role key is server-only and
never reaches the browser. \\
\bottomrule
\end{longtable}

\begin{lstlisting}[language=TS,caption={lib/supabase/public.ts --- read the comment; it is the whole justification}]
// Stateless, cookie-free Supabase client using the public anon key. RLS still
// applies (anon can only read approved hotels and other public rows). Safe to
// use inside `unstable_cache`, which must not touch request-bound data like
// cookies -- unlike the SSR `server.ts` client.
export function createPublicClient() {
  return createClient(
    process.env.NEXT_PUBLIC_SUPABASE_URL!,
    process.env.NEXT_PUBLIC_SUPABASE_ANON_KEY!,
    { auth: { persistSession: false, autoRefreshToken: false } },
  );
}
\end{lstlisting}

\saythis{Four clients, four reasons. The browser one manages cookies. The server
one reads the request's cookies so RLS knows who you are. The public one is
cookie-free because Next.js caching forbids request-scoped data inside a cached
function. And the admin one uses the service-role key to bypass RLS --- we use it
in exactly one place, the Razorpay webhook, because there is no user session
when a payment gateway calls us.}

\subsection{Row Level Security, properly explained}

Normally, an application checks permissions:
\emph{``is this user allowed to see this booking? ok, run the query.''}
If you forget that check in one endpoint, the data leaks.

RLS inverts it. You attach a rule to the \emph{table}, and Postgres appends it to
every single query, forever, no matter who is asking or from where.

\begin{lstlisting}[language=SQL,caption={supabase/booking\_schema.sql}]
alter table public.bookings enable row level security;

create policy "bookings: read own"
  on public.bookings for select using (auth.uid() = guest_id);

create policy "bookings: read manager"
  on public.bookings for select
  using (exists (
    select 1 from public.hotels h
    where h.id = hotel_id and h.manager_id = auth.uid()
  ));

create policy "bookings: read admin"
  on public.bookings for select using (public.get_my_role() = 'admin');
\end{lstlisting}

Now \texttt{select * from bookings} returns:
your own bookings if you are a guest;
your hotels' bookings if you are a manager;
everything if you are an admin;
and \textbf{nothing} if you are a stranger --- even with a valid API key.
Multiple \texttt{SELECT} policies are OR-ed together.

\texttt{auth.uid()} is a Postgres function Supabase provides. It reads the user's
ID out of the JWT attached to the request.

\subsection{The RLS recursion problem (and how we solved it)}

This is a genuinely hard problem we hit, and explaining it will impress anyone.

Consider: the policy on \texttt{hotels} needs to check ``is the caller assigned
staff on this hotel?'', which requires reading \texttt{hotel\_staff}. But the
policy on \texttt{hotel\_staff} needs to check ``does the caller own this
hotel?'', which requires reading \texttt{hotels}. Each table's policy queries the
other table --- whose policy queries the first table. Postgres detects the loop
and throws \texttt{infinite recursion detected in policy}.

\textbf{The fix} is \texttt{SECURITY DEFINER}. A function declared
\texttt{SECURITY DEFINER} runs with the privileges of the user who \emph{created}
it (the database owner), not the user calling it --- so RLS does not apply to the
queries \emph{inside} it. It reads one table, cleanly, and returns a boolean.

\begin{lstlisting}[language=SQL,caption={supabase/hms\_management\_migration.sql}]
-- Single-table definer helpers used inside RLS policies to AVOID the mutual
-- recursion hotels <-> hotel_staff. These bypass RLS internally.

create or replace function public.owns_hotel(p_hotel_id uuid)
returns boolean language sql stable
security definer set search_path = public
as $$ select exists (select 1 from public.hotels h
                     where h.id = p_hotel_id and h.manager_id = auth.uid()); $$;

-- True if the caller is the hotel's manager OR staff holding `p_perm`.
create or replace function public.can_manage_hotel(p_hotel_id uuid, p_perm text)
returns boolean language sql stable
security definer set search_path = public
as $$
  select
    exists (select 1 from public.hotels h
            where h.id = p_hotel_id and h.manager_id = auth.uid())
    or exists (select 1 from public.hotel_staff s
               where s.hotel_id = p_hotel_id and s.staff_id = auth.uid()
                 and p_perm = any (s.permissions));
$$;
\end{lstlisting}

\begin{warnbox}
\textbf{Why \texttt{set search\_path = public}?} A \texttt{SECURITY DEFINER}
function runs as a privileged user. If an attacker could change the schema search
path, they could make the function call \emph{their} malicious
\texttt{public.hotels} instead of ours. Pinning \texttt{search\_path} closes that
hole. Every definer function in this project pins it. This is the single most
commonly forgotten line in Postgres security, and we did not forget it.
\end{warnbox}

The same technique is used in messaging:
\texttt{can\_access\_conversation()} and \texttt{validate\_message\_sender()}
exist purely so the \texttt{messages} policies never have to query
\texttt{conversations} under RLS.

\subsection{The trigger that creates every profile}

Supabase stores users in \texttt{auth.users}, a schema we do not control. We need
our own \texttt{public.profiles} row for each of them (holding the role). We do
this with a trigger, so it is impossible for a user to exist without a profile.

\begin{lstlisting}[language=SQL,caption={supabase/schema.sql}]
create or replace function public.handle_new_user()
returns trigger language plpgsql
security definer set search_path = public
as $$
begin
  insert into public.profiles (id, role, full_name, phone, dob, location)
  values (
    new.id,
    coalesce(new.raw_user_meta_data ->> 'role', 'guest'),  -- default: guest
    new.raw_user_meta_data ->> 'full_name',
    new.raw_user_meta_data ->> 'phone',
    nullif(new.raw_user_meta_data ->> 'dob', '')::date,
    new.raw_user_meta_data ->> 'location'
  )
  on conflict (id) do nothing;

  -- Managers additionally get a pending verification row, automatically.
  if coalesce(new.raw_user_meta_data ->> 'role', 'guest') = 'manager' then
    insert into public.manager_verifications (
      user_id, business_name, registration_number, business_address,
      document_url, status)
    values (new.id,
      new.raw_user_meta_data ->> 'business_name',
      new.raw_user_meta_data ->> 'registration_number',
      new.raw_user_meta_data ->> 'business_address',
      new.raw_user_meta_data ->> 'document_url',
      'pending');
  end if;
  return new;
end; $$;

create trigger on_auth_user_created
  after insert on auth.users
  for each row execute function public.handle_new_user();
\end{lstlisting}

\texttt{raw\_user\_meta\_data} is the JSON blob we passed in the
\texttt{options.data} of \texttt{supabase.auth.signUp()}. So the client says
``I am signing up as a manager, here is my business name'', and the database
turns that into a profile \emph{plus} a pending verification request --- in one
transaction, before the user has even confirmed their email.

\saythis{We never insert a profile from application code. A Postgres trigger on
\texttt{auth.users} does it, reading the sign-up metadata. That means a user row
without a matching profile row is structurally impossible --- there is no code
path that can forget.}

\subsection{Why so much logic lives in the database (RPC functions)}

A large amount of this project's business logic is written in PL/pgSQL and called
via \texttt{supabase.rpc("function\_name", \{...\})}. That is a deliberate
architectural choice, and here is the reasoning.

\begin{enumerate}[leftmargin=*]
  \item \textbf{Atomicity.} \texttt{book\_room} reserves inventory for every night
        \emph{and} writes the booking \emph{and} writes the payment row. If any
        step fails, Postgres rolls the whole thing back. Doing this from
        JavaScript across three separate HTTP calls would leave orphaned data
        the moment the network hiccups.
  \item \textbf{Concurrency.} Only the database can take a row lock. See
        Section~\ref{sec:concurrency}.
  \item \textbf{Trust.} The price is computed in SQL from the room's real price.
        A malicious client that POSTs \texttt{\{"total": 1\}} is ignored --- the
        server never reads a price from the request body.
  \item \textbf{One place to change.} Both the web app and the webhook confirm
        payments. They both call the same function.
\end{enumerate}

\subsection{The complete table catalogue}

\begin{longtable}{@{}p{3.5cm}p{11.3cm}@{}}
\toprule
\textbf{Table} & \textbf{Purpose and key columns} \\
\midrule
\endhead
\texttt{profiles} & One row per auth user. Holds \texttt{role}
(guest / manager / staff / admin), name, phone, \texttt{suspended}. \\
\addlinespace
\texttt{manager\_verifications} & A manager's business documents and approval
\texttt{status} (pending / approved / rejected / more\_info). \\
\addlinespace
\texttt{hotels} & The listing. \texttt{status} gates visibility
(draft $\to$ pending $\to$ approved). Also holds lat/lng, amenities array,
policies, \texttt{gst\_percent}, \texttt{deleted\_at} (soft delete). \\
\addlinespace
\texttt{rooms} & A room \emph{type} (not a physical room). \texttt{price},
\texttt{capacity}, \texttt{total\_units} = how many of this type exist. \\
\addlinespace
\texttt{room\_inventory} & \textbf{The heart of availability.} One row per
(room type, date), holding \texttt{total\_count} and \texttt{booked\_count},
with a CHECK constraint that booked never exceeds total. \\
\addlinespace
\texttt{bookings} & \texttt{status} in pending / confirmed / checked\_in /
completed / cancelled. Snapshots the price at booking time. \texttt{guest\_id}
is nullable to support walk-in bookings. \\
\addlinespace
\texttt{payments} & One per booking. \texttt{order\_id} (Razorpay order),
\texttt{transaction\_id} (Razorpay payment id), \texttt{status}. \\
\addlinespace
\texttt{conversations} & One per (hotel, guest) pair --- enforced by a UNIQUE
constraint. Holds \texttt{guest\_unread}, \texttt{host\_unread}, and a preview
of the last message. \\
\addlinespace
\texttt{messages} & \texttt{body}, \texttt{attachments} (jsonb array),
\texttt{sender\_role}. \\
\addlinespace
\texttt{hotel\_staff} & Assignment of a staff account to one hotel, with a
\texttt{permissions text[]} array. \\
\addlinespace
\texttt{staff\_invites} & Pending email invitations, with a random
\texttt{token} and a 14-day \texttt{expires\_at}. \\
\addlinespace
\texttt{reviews} & rating 1--5 + comment. \\
\addlinespace
\texttt{hotel\_photos}, \texttt{room\_photos} & Ordered image URLs. \\
\addlinespace
\texttt{admin\_audit\_log} & Every admin action, with a jsonb \texttt{details}
blob. Append-only; admins can read but the app never deletes. \\
\bottomrule
\end{longtable}

\subsection{Database functions, at a glance}

\begin{longtable}{@{}p{5.2cm}p{9.6cm}@{}}
\toprule
\textbf{Function} & \textbf{What it does} \\
\midrule
\endhead
\texttt{get\_my\_role()} & Returns the caller's role without tripping RLS
recursion on \texttt{profiles}. Used inside dozens of policies. \\
\addlinespace
\texttt{handle\_new\_user()} & Trigger. Creates the profile on signup. \\
\addlinespace
\texttt{room\_availability(hotel, in, out)} & Returns each room type with the
number of units free for the \emph{whole} range. \\
\addlinespace
\texttt{book\_room(...)} & Atomically reserves inventory, creates the booking and
the pending payment. Returns the booking id. \\
\addlinespace
\texttt{book\_offline(...)} & Same, for a walk-in guest with no account. \\
\addlinespace
\texttt{attach\_payment\_order(booking, order)} & Stores the Razorpay order id. \\
\addlinespace
\texttt{confirm\_razorpay\_payment(...)} & Marks payment completed + booking
confirmed. Only after the server has verified the signature. \\
\addlinespace
\texttt{cancel\_booking(booking, reason)} & Releases inventory night by night and
computes the refund from the time-based policy. \\
\addlinespace
\texttt{check\_in\_booking(id)} / \texttt{check\_out\_booking(id)} & Front-desk
state transitions, guarded by \texttt{can\_operate\_hotel()}. \\
\addlinespace
\texttt{can\_manage\_hotel(hotel, perm)} & Manager OR staff-with-permission. \\
\addlinespace
\texttt{can\_operate\_hotel(hotel)} & Manager OR any assigned staff. \\
\addlinespace
\texttt{shares\_conversation\_with(profile)} & Powers the privacy rule that lets
a host read a guest's profile only when a conversation exists. \\
\addlinespace
\texttt{get\_hotel\_reviews(hotel)} & Joins reviews to profiles internally and
returns \emph{only} safe fields (name, join date) --- never phone or email. \\
\addlinespace
\texttt{admin\_review\_hotel(...)}, \texttt{admin\_review\_manager(...)},
\texttt{admin\_set\_user\_role(...)}, \texttt{admin\_set\_user\_suspended(...)}
& Every one begins with
\texttt{if get\_my\_role() <> 'admin' then raise exception} and ends by writing
to \texttt{admin\_audit\_log}. \\
\bottomrule
\end{longtable}

\subsection{Storage buckets}

\begin{longtable}{@{}p{4.2cm}p{2.0cm}p{8.6cm}@{}}
\toprule
\textbf{Bucket} & \textbf{Public?} & \textbf{Policy} \\
\midrule
\endhead
\texttt{manager-documents} & No & Insert allowed anonymously (the upload happens
\emph{during} signup, before a session exists); read restricted to admins. \\
\addlinespace
\texttt{hotel-images} & Yes & Anyone reads; authenticated users write. \\
\addlinespace
\texttt{message-attachments} & No & Made private in a later migration. Reads go
through a signed-URL endpoint that checks conversation membership. \\
\bottomrule
\end{longtable}

\subsection{Performance: the indexes that matter}

\begin{lstlisting}[language=SQL,caption={supabase/perf\_migration.sql}]
-- Trigram matching lets `ILIKE '%term%'` (leading-wildcard) use an index
-- instead of scanning every row. Required for instant text search.
create extension if not exists pg_trgm;

create index if not exists hotels_location_trgm_idx
  on public.hotels using gin (location gin_trgm_ops);
create index if not exists hotels_name_trgm_idx
  on public.hotels using gin (name gin_trgm_ops);

-- The catalog query is: status = 'approved' ORDER BY created_at DESC.
-- A partial index over ONLY approved rows serves the filter and the sort
-- with no extra sort step.
create index if not exists hotels_approved_created_idx
  on public.hotels (created_at desc)
  where status = 'approved';
\end{lstlisting}

\saythis{Searching for ``goa'' means \texttt{ILIKE '\%goa\%'}, and a normal B-tree
index cannot help with a leading wildcard --- Postgres would scan every hotel. So
we enabled the \texttt{pg\_trgm} extension and built GIN trigram indexes, which
index every three-character substring. And because the catalogue query always
filters on \texttt{status = 'approved'}, we used a \emph{partial} index --- it only
contains approved rows, so it is smaller and it satisfies the sort as well as the
filter.}

\newpage
%==============================================================================
\section{Authentication}

\subsection{The three sign-up paths}

\subsubsection{Guest, with email and password}

\begin{lstlisting}[language=TS,caption={lib/auth.ts}]
export async function signUpGuest(input: GuestSignupInput) {
  const supabase = createClient();
  const { data, error } = await supabase.auth.signUp({
    email: input.email,
    password: input.password,
    options: {
      emailRedirectTo: `${siteUrl}/auth/callback`,
      data: {                         // -> becomes raw_user_meta_data
        role: "guest",
        full_name: input.fullName,
        dob: input.dob,
        phone: input.phone,
        location: input.location,
      },
    },
  });
  if (error) throw error;
  return data;
}
\end{lstlisting}

We never write to \texttt{profiles} here. We hand the data to Supabase Auth, and
the \texttt{handle\_new\_user} trigger creates the profile. This works even
before the user confirms their email.

\subsubsection{Manager, with a verification document}

The document must be uploaded \emph{before} \texttt{signUp()} returns, because
the trigger needs its path in the metadata. This is why the
\texttt{manager-documents} bucket permits anonymous inserts.

\begin{lstlisting}[language=TS,caption={lib/auth.ts}]
const path = `pending/${crypto.randomUUID()}.${ext}`;
await supabase.storage.from("manager-documents").upload(path, input.document);

await supabase.auth.signUp({
  email: input.email, password: input.password,
  options: { data: { role: "manager", business_name: ..., document_url: path } },
});
\end{lstlisting}

\subsubsection{Google OAuth}

\begin{lstlisting}[language=TS,caption={lib/auth.ts --- note the origin trick}]
export async function signInWithGoogle() {
  const supabase = createClient();
  // Always runs in the browser (click handler), so prefer the live origin.
  // This keeps localhost, Vercel previews, and production each redirecting
  // back to themselves rather than a single hard-coded site URL.
  const origin = typeof window !== "undefined" ? window.location.origin : siteUrl;
  return supabase.auth.signInWithOAuth({
    provider: "google",
    options: { redirectTo: `${origin}/auth/callback` },
  });
}
\end{lstlisting}

Google only gives us an email and a name. It gives us no phone number and no date
of birth. So the trigger creates a profile with those fields null, and
\file{app/auth/callback/route.ts} notices and sends the user to
\texttt{/onboarding} to fill them in.

\subsection{The callback route: one door, four destinations}

\begin{lstlisting}[language=TS,caption={app/auth/callback/route.ts}]
export async function GET(request: Request) {
  const code = searchParams.get("code");
  if (!code) return NextResponse.redirect(`${origin}/login?error=missing_code`);

  const supabase = await createClient();
  const { error } = await supabase.auth.exchangeCodeForSession(code);
  if (error) return NextResponse.redirect(`${origin}/login?error=auth_failed`);

  const { data: profile } = await supabase
    .from("profiles").select("role, phone").eq("id", user.id).single();
  const role = profile?.role ?? "guest";

  if (role === "admin")   return redirect(`${origin}/admin/dashboard`);

  if (role === "manager") {                       // approved? -> dashboard
    const { data: mv } = await supabase           // else      -> waiting room
      .from("manager_verifications").select("status")
      .eq("user_id", user.id).order("created_at", { ascending: false })
      .limit(1).maybeSingle();
    return redirect(`${origin}${mv?.status === "approved"
      ? "/manager/dashboard" : "/manager/waiting"}`);
  }

  // Guest. OAuth sign-ups land here with no phone -> finish the profile first.
  if (!profile?.phone) return redirect(`${origin}/onboarding`);
  return redirect(`${origin}/`);
}
\end{lstlisting}

Both the email-confirmation link and the Google redirect land here. The
\texttt{code} is a one-time authorisation code; \texttt{exchangeCodeForSession}
swaps it for a session cookie. This is the PKCE flow, and Supabase implements it
for us.

\subsection{\texttt{getUser} vs \texttt{getClaims} vs \texttt{getSession}}

We use all three, in different places, on purpose. If you can explain this, you
understand the auth system better than most working developers.

\begin{longtable}{@{}p{2.6cm}p{4.4cm}p{7.8cm}@{}}
\toprule
\textbf{Call} & \textbf{Cost} & \textbf{Where we use it, and why} \\
\midrule
\endhead
\texttt{getSession()} & Free (reads local storage) & \file{components/AuthProvider.tsx}.
It is only driving UI state (``show the avatar''). It can be trusted for
cosmetics, never for authorisation. \\
\addlinespace
\texttt{getClaims()} & Cheap (verifies the JWT signature locally) &
\file{lib/supabase/middleware.ts} and \file{lib/authServer.ts}. We only need the
user id, and a forged token fails signature verification. \\
\addlinespace
\texttt{getUser()} & Expensive (network call to Supabase Auth) &
Route Handlers that write --- \file{api/bookings}, \file{api/payments/*}. Here we
want the strongest possible guarantee before touching money. \\
\bottomrule
\end{longtable}

\begin{lstlisting}[language=TS,caption={components/AuthProvider.tsx --- the comment explains the trade-off}]
// getSession() reads the token from local storage (instant, no network
// round-trip), unlike getUser() which validates against Supabase's servers
// on every page. For client-side UI state this is the recommended call;
// anything security-sensitive is enforced server-side / by RLS anyway.
const { data: { session } } = await supabase.auth.getSession();
\end{lstlisting}

\saythis{\texttt{getSession} just reads the token from local storage --- fast, but a
user could edit local storage, so we only use it to decide what to \emph{draw}.
\texttt{getClaims} cryptographically verifies the token's signature locally, which
is enough to trust the user id without a network call --- that is what middleware
uses. \texttt{getUser} actually asks Supabase's server, and we save it for the
routes that create bookings and payments.}

\subsection{How the session survives}

\begin{enumerate}[leftmargin=*,noitemsep]
  \item Supabase issues a short-lived JWT (access token) and a refresh token.
  \item \texttt{@supabase/ssr} stores them in \textbf{cookies}, not just local
        storage, so the \emph{server} can read them too.
  \item On a protected route, \file{proxy.ts} runs, and the act of constructing
        the server client refreshes the token if needed and writes the new cookie
        onto the response.
  \item On public routes we skip middleware entirely --- the browser client
        refreshes its own token in the background.
\end{enumerate}

That fourth point is why the \texttt{matcher} in \file{proxy.ts} is safe to
restrict.

\newpage
%==============================================================================
\section{The Booking System}

\subsection{The model: room \emph{types}, not room numbers}

We do not store ``Room 302''. We store a room \emph{type} --- ``Deluxe Suite'' ---
with a \texttt{total\_units} count of how many exist. Availability is therefore a
counting problem, not an assignment problem.

\begin{notebox}
\textbf{Why?} Because guests do not book Room 302; they book ``a deluxe room''.
The hotel assigns the physical room at check-in. Modelling it as a count is both
simpler and closer to how hotels actually operate. It also means the front desk
does not have to maintain a floor plan in our system.
\end{notebox}

\subsection{\texttt{room\_inventory}: one row per room type per night}

\begin{lstlisting}[language=SQL]
create table public.room_inventory (
  id           uuid primary key default gen_random_uuid(),
  room_id      uuid not null references public.rooms (id) on delete cascade,
  date         date not null,
  total_count  int not null,
  booked_count int not null default 0,
  unique (room_id, date),
  constraint inventory_not_overbooked check (booked_count <= total_count)
);
\end{lstlisting}

Two details worth pointing at:

\begin{itemize}[leftmargin=*,noitemsep]
  \item \texttt{unique (room\_id, date)} --- there can only ever be one inventory
        row per night. This is what makes the lock in \texttt{book\_room} work.
  \item \texttt{check (booked\_count <= total\_count)} --- a database-level
        guarantee that overbooking is \emph{impossible}, even if every line of
        application code were wrong.
\end{itemize}

Rows are created \textbf{lazily}. A hotel with 20 room types open for two years
would otherwise need 14,600 rows on day one. Instead the row for a given night
springs into existence the first time somebody books that night.

\subsection{Checking availability}

\begin{lstlisting}[language=SQL,caption={supabase/booking\_schema.sql}]
-- A room is bookable for the whole range only if it has stock on EVERY night,
-- so availability = total_units - max(booked nights in range).
create or replace function public.room_availability(
  p_hotel_id uuid, p_check_in date, p_check_out date)
returns table (room_id uuid, name text, price numeric,
               capacity int, amenities text[], available int)
language sql stable security definer set search_path = public
as $$
  select r.id, r.name, r.price, r.capacity, r.amenities,
    r.total_units - coalesce((
      select max(i.booked_count)
      from public.room_inventory i
      where i.room_id = r.id
        and i.date >= p_check_in
        and i.date <  p_check_out       -- note: check-out night is NOT occupied
    ), 0) as available
  from public.rooms r
  where r.hotel_id = p_hotel_id
  order by r.price asc;
$$;
\end{lstlisting}

The clever bit is \texttt{max(booked\_count)}. If a guest wants three nights and
the room is fully booked on the middle night, they cannot have the range at all.
Taking the maximum booked count across the range gives exactly the number of
units free for the \emph{entire} stay.

\texttt{i.date < p\_check\_out} (not \texttt{<=}) encodes the fact that you do not
occupy a room on the night of your check-out day.

\subsection{Creating the booking, atomically}
\label{sec:concurrency}

\textbf{This is the most important 30 lines in the project.} Learn them.

\begin{lstlisting}[language=SQL,caption={supabase/booking\_schema.sql --- book\_room}]
create or replace function public.book_room(
  p_room_id uuid, p_check_in date, p_check_out date,
  p_guests int, p_num_rooms int, p_special text)
returns uuid language plpgsql security definer set search_path = public
as $$
declare v_room public.rooms; v_hotel public.hotels;
        v_uid uuid := auth.uid(); v_nights int; v_avail int; d date;
        v_base numeric; v_gst numeric; v_service_charge numeric;
        v_fee numeric; v_total numeric; v_id uuid;
begin
  if v_uid is null then raise exception 'Not authenticated'; end if;
  if p_check_out <= p_check_in then raise exception 'Check-out must be after check-in'; end if;
  if p_check_in < current_date then raise exception 'Check-in cannot be in the past'; end if;
  if p_num_rooms < 1 then raise exception 'At least one room is required'; end if;

  select * into v_room  from public.rooms  where id = p_room_id;
  select * into v_hotel from public.hotels where id = v_room.hotel_id;
  v_nights := p_check_out - p_check_in;

  -- Reserve each night: lazily create the inventory row, LOCK it, then book.
  d := p_check_in;
  while d < p_check_out loop
    insert into public.room_inventory (room_id, date, total_count, booked_count)
      values (p_room_id, d, v_room.total_units, 0)
      on conflict (room_id, date) do nothing;         -- (1) idempotent create

    select total_count - booked_count into v_avail
      from public.room_inventory
      where room_id = p_room_id and date = d
      for update;                                     -- (2) THE ROW LOCK

    if v_avail < p_num_rooms then
      raise exception 'Only % room(s) left on %', v_avail, to_char(d, 'Mon DD');
    end if;                                           -- (3) rolls back everything

    update public.room_inventory
      set booked_count = booked_count + p_num_rooms
      where room_id = p_room_id and date = d;

    d := d + 1;
  end loop;

  -- Price is computed HERE, from the DB's own numbers. Never from the client.
  v_base           := v_room.price * p_num_rooms * v_nights;
  v_gst            := round(v_base * (coalesce(v_hotel.gst_percent, 18.00) / 100.0), 2);
  v_service_charge := round(v_base * (coalesce(v_hotel.service_charge_percent, 0) / 100.0), 2);
  v_fee            := round(v_base * 0.02, 2);        -- platform fee
  v_total          := v_base + v_gst + v_service_charge + v_fee;

  insert into public.bookings (...) values (..., 'pending', p_special)
    returning id into v_id;
  insert into public.payments (booking_id, amount, payment_method, status)
    values (v_id, v_total, 'upi', 'pending');

  return v_id;
end; $$;
\end{lstlisting}

\subsubsection{Why \texttt{FOR UPDATE} is the whole ballgame}

Imagine one room left, and two guests click ``Book'' in the same millisecond.

\begin{lstlisting}
WITHOUT the lock:
  Alice reads booked_count = 4, total = 5  -> "1 free, ok!"
  Bob   reads booked_count = 4, total = 5  -> "1 free, ok!"
  Alice writes booked_count = 5
  Bob   writes booked_count = 5            <-- Bob overwrote Alice. OVERBOOKED.

WITH `for update`:
  Alice's SELECT ... FOR UPDATE locks the row for date D.
  Bob's SELECT ... FOR UPDATE *blocks* -- it waits.
  Alice writes booked_count = 5, her transaction commits, lock released.
  Bob's SELECT now returns booked_count = 5 -> v_avail = 0
  Bob raises 'Only 0 room(s) left on Mar 14' -> his whole transaction ROLLS BACK.
\end{lstlisting}

And because the entire function is one transaction, when Bob's exception fires on
night three, the reservations he already made for nights one and two are undone
automatically. There is no cleanup code, and there cannot be a partial booking.

\saythis{Two people booking the last room at the same instant is a classic race
condition. We solve it the way a bank does: \texttt{SELECT ... FOR UPDATE} takes a
row-level lock on each night's inventory row, so the second transaction physically
waits for the first to commit and then sees the updated count. The whole function
is one transaction, so if night three is full, nights one and two are rolled back
automatically. And there is a CHECK constraint as a last line of defence --- even a
bug could not overbook the hotel.}

\subsection{Why the price is calculated twice}

\file{lib/booking.ts} has \texttt{computeQuote()}, which mirrors the SQL exactly.
This is not duplication by accident:

\begin{lstlisting}[language=TS,caption={lib/booking.ts}]
export const GST_RATE = 0.18;
export const PLATFORM_FEE_RATE = 0.02;

// Mirror of the server-side price math in book_room (kept in sync for preview).
export function computeQuote(roomPrice, nights, numRooms,
                             gstPercent = 18, serviceChargePercent = 0): PriceQuote {
  const base          = round2(roomPrice * nights * numRooms);
  const gst           = round2(base * (gstPercent / 100));
  const serviceCharge = round2(base * (serviceChargePercent / 100));
  const platformFee   = round2(base * PLATFORM_FEE_RATE);
  const total = round2(base + gst + serviceCharge + platformFee);
  return { roomPrice, nights, numRooms, base, gst, serviceCharge, platformFee, total };
}
\end{lstlisting}

The TypeScript version exists \emph{only} to render the price breakdown while the
user is choosing dates, with no network call. The number it produces is never
sent to the server and never trusted. \texttt{book\_room} recomputes from
\texttt{rooms.price} and \texttt{hotels.gst\_percent}.

\saythis{The client-side calculation is a preview, for instant feedback. The
authoritative number comes from \texttt{book\_room}, which reads the price out of
the database. If you tampered with the JavaScript to say the room costs one rupee,
the booking would still be created for the real amount --- the API never accepts a
price as input.}

\subsection{Cancellation and the refund policy}

\begin{lstlisting}[language=SQL,caption={supabase/booking\_schema.sql --- cancel\_booking}]
  v_hours := extract(epoch from (v_b.check_in::timestamp - now())) / 3600;
  if    v_hours >= 48 then v_pct := 1.0;   -- full refund
  elsif v_hours >= 24 then v_pct := 0.5;   -- half
  else                     v_pct := 0.0;   -- none
  end if;
  v_refund := round(v_b.total_price * v_pct, 2);

  -- Release the reserved inventory, night by night.
  d := v_b.check_in;
  while d < v_b.check_out loop
    update public.room_inventory
      set booked_count = greatest(0, booked_count - v_b.num_rooms)
      where room_id = v_b.room_id and date = d;
    d := d + 1;
  end loop;
\end{lstlisting}

\texttt{greatest(0, ...)} is defensive: even if something else had already
decremented the count, we can never drive it negative.

There is a mirror of this policy in TypeScript too (\texttt{refundPolicy()} in
\file{lib/booking.ts}) --- again, display only. The comment in the source says so
explicitly: \emph{``Refund policy mirror (display only --- server recomputes
authoritatively).''}

\subsection{Who is allowed to cancel}

\begin{lstlisting}[language=SQL]
if v_b.guest_id <> auth.uid() and public.get_my_role() <> 'admin' then
  raise exception 'Not allowed';
end if;
\end{lstlisting}

Later, \file{hms\_management\_migration.sql} \emph{redefines} the same function to
also permit the hotel's manager and staff, so the front desk can cancel a booking
for a guest who calls the hotel. This is why the same function name appears in two
migration files --- the later one wins. That is intentional, and it is how
migrations work: they are applied in order, and each is idempotent
(\texttt{create or replace}, \texttt{if not exists}).

\newpage
%==============================================================================
\section{Payments: Razorpay}

\subsection{Why a gateway at all}

We never see a card number. If we did, we would be legally responsible for PCI-DSS
compliance. Razorpay's checkout runs on \emph{Razorpay's} domain, inside a modal.
The card details go from the user's browser straight to Razorpay. Our server only
ever handles opaque identifiers.

\subsection{The three-step handshake}

\begin{lstlisting}
BROWSER                      OUR SERVER                    RAZORPAY
   |                              |                            |
   |-- POST /api/payments/------->|                            |
   |     razorpay/order           |-- orders.create(amount) -->|
   |     { bookingId }            |<---------- order_id -------|
   |                              |                            |
   |                              |  attach_payment_order()    |
   |<---- { orderId, keyId } -----|  (store order_id in DB)    |
   |                                                           |
   |------ open checkout.js modal, user pays ----------------->|
   |<----- { payment_id, order_id, signature } ----------------|
   |                              |                            |
   |-- POST /api/payments/------->|                            |
   |     razorpay/verify          |  HMAC_SHA256(              |
   |                              |    order_id + "|" + payment_id,
   |                              |    RAZORPAY_KEY_SECRET )   |
   |                              |  == signature ?            |
   |                              |                            |
   |                              |  confirm_razorpay_payment()|
   |<-------- { ok: true } -------|  booking -> 'confirmed'    |
\end{lstlisting}

\subsection{Step 1: creating the order (server)}

\begin{lstlisting}[language=TS,caption={app/api/payments/razorpay/order/route.ts (abridged)}]
const { data: { user } } = await supabase.auth.getUser();
if (!user) return NextResponse.json({ error: "Not authenticated" }, { status: 401 });

// RLS ensures the booking belongs to this user.
const { data: booking } = await supabase
  .from("bookings")
  .select("total_price, base_price, status, hotel_id, hotels (require_advance, ...)")
  .eq("id", bookingId).maybeSingle();

if (!booking) return NextResponse.json({ error: "Booking not found" }, { status: 404 });
if (booking.status !== "pending")
  return NextResponse.json({ error: "This booking is not awaiting payment" }, { status: 400 });

// Prices are stored in INR, so the order amount is simply paise.
const amountPaise = Math.round(amountToCharge * 100);

const order = await razorpay.orders.create({
  amount: amountPaise, currency: "INR",
  receipt: bookingId, notes: { bookingId },
  ...(transfers ? { transfers } : {}),   // Razorpay Route split, see below
});

await supabase.rpc("attach_payment_order", {
  p_booking_id: bookingId, p_order_id: order.id,
});

return NextResponse.json({
  orderId: order.id, amount: amountPaise, currency: "INR",
  keyId: process.env.NEXT_PUBLIC_RAZORPAY_KEY_ID,   // publishable, safe
});
\end{lstlisting}

Three deliberate details:

\begin{itemize}[leftmargin=*,noitemsep]
  \item The amount comes from \texttt{booking.total\_price} in the database. The
        request body contains \emph{only} a \texttt{bookingId}.
  \item ``Booking not found'' is what an attacker sees when they pass someone
        else's booking id --- because RLS returned zero rows. We did not have to
        write an ownership check; the database performed it.
  \item Razorpay works in \textbf{paise}, the smallest currency unit. Rs.~1{,}500
        is \texttt{150000}. Using integers avoids floating-point money bugs.
\end{itemize}

\subsection{Step 2: the browser modal}

\begin{lstlisting}[language=TS,caption={lib/razorpayCheckout.ts}]
const ok = await loadScript("https://checkout.razorpay.com/v1/checkout.js");

const rzp = new window.Razorpay({
  key: order.keyId, order_id: order.orderId,
  amount: order.amount, currency: order.currency,
  name: "BookNest",
  handler: async (resp: RazorpayResponse) => {         // Razorpay calls this on success
    const v = await fetch("/api/payments/razorpay/verify", {
      method: "POST",
      body: JSON.stringify({
        bookingId,
        orderId:   resp.razorpay_order_id,
        paymentId: resp.razorpay_payment_id,
        signature: resp.razorpay_signature,
      }),
    });
    if (v.ok) onSuccess(); else onError(...);
  },
  modal: { ondismiss: () => onError("Payment was cancelled.") },
});
rzp.open();
\end{lstlisting}

\subsection{Step 3: verifying the signature (server)}

\begin{lstlisting}[language=TS,caption={app/api/payments/razorpay/verify/route.ts}]
// Signature = HMAC_SHA256(order_id + "|" + payment_id, RAZORPAY_KEY_SECRET)
const expected = crypto
  .createHmac("sha256", process.env.RAZORPAY_KEY_SECRET!)
  .update(`${orderId}|${paymentId}`)
  .digest("hex");

// Timing-safe comparison.
const valid =
  expected.length === signature.length &&
  crypto.timingSafeEqual(Buffer.from(expected), Buffer.from(signature));

if (!valid)
  return NextResponse.json({ error: "Payment signature verification failed" }, { status: 400 });

const { error } = await supabase.rpc("confirm_razorpay_payment", {
  p_booking_id: bookingId, p_order_id: orderId, p_payment_id: paymentId,
});
\end{lstlisting}

\begin{notebox}
\textbf{What is an HMAC, in plain English?} It is a fingerprint of a message,
made using a secret only two parties know. Razorpay knows our
\texttt{KEY\_SECRET}. We know our \texttt{KEY\_SECRET}. Nobody else does. So when
Razorpay sends back a fingerprint of \texttt{order\_id|payment\_id} and we compute
the same fingerprint and it matches, we know the message really came from Razorpay
and was not altered. A hacker cannot forge it without the secret.
\end{notebox}

\begin{notebox}
\textbf{Why \texttt{timingSafeEqual} and not \texttt{===}?} Because \texttt{===}
on strings stops comparing at the first mismatched character. An attacker who can
measure response times very precisely could learn the signature one character at a
time --- a \emph{timing attack}. \texttt{crypto.timingSafeEqual} always compares
every byte, so the comparison takes constant time regardless of where the strings
differ.
\end{notebox}

\saythis{Anyone can POST to our verify endpoint with a made-up payment id. The
signature is what makes it trustworthy: it is an HMAC-SHA256 of the order id and
payment id, keyed with a secret that only Razorpay and our server hold. We
recompute it and compare in constant time. If it does not match, nothing happens.}

\subsection{Three roads to ``confirmed'' --- and why we need all three}

A payment can succeed while the user's Wi-Fi dies, or their phone battery goes
flat, or they close the tab. The money left their account. Our database still says
\texttt{pending}. That is unacceptable, so there are three independent paths:

\begin{longtable}{@{}p{2.5cm}p{4.1cm}p{8.2cm}@{}}
\toprule
\textbf{Path} & \textbf{Trigger} & \textbf{Mechanism} \\
\midrule
\endhead
\textbf{Happy} & Razorpay's JS calls our \texttt{handler} &
\file{api/payments/razorpay/verify} verifies the signature and confirms. \\
\addlinespace
\textbf{Webhook} & Razorpay's \emph{server} POSTs to us &
\file{api/webhooks/razorpay} verifies an HMAC of the raw request body against
\texttt{RAZORPAY\_WEBHOOK\_SECRET}, then confirms using the service-role client
(there is no user session). \\
\addlinespace
\textbf{Reconcile} & The user re-opens the booking page &
\file{lib/reconcile.ts} asks Razorpay ``was order X actually paid?'' via
\texttt{orders.fetchPayments()} and confirms if a captured payment exists. \\
\bottomrule
\end{longtable}

\begin{lstlisting}[language=TS,caption={lib/reconcile.ts --- the webhook-free safety net}]
export async function reconcilePendingBooking(supabase, bookingId, orderId) {
  if (!orderId) return false;
  try {
    const res = await razorpay.orders.fetchPayments(orderId);
    const captured = res.items?.find((p) => p.status === "captured");
    if (!captured) return false;

    const { error } = await supabase.rpc("confirm_razorpay_payment", {
      p_booking_id: bookingId, p_order_id: orderId, p_payment_id: captured.id,
    });
    return !error;
  } catch { return false; }
}
\end{lstlisting}

\begin{lstlisting}[language=TS,caption={app/api/webhooks/razorpay/route.ts --- signature over the RAW body}]
const signature = request.headers.get("x-razorpay-signature") ?? "";
const raw = await request.text();          // must be the raw text, not parsed JSON

const expected = crypto.createHmac("sha256", secret).update(raw).digest("hex");
const valid = expected.length === signature.length &&
  crypto.timingSafeEqual(Buffer.from(expected), Buffer.from(signature));
if (!valid) return NextResponse.json({ error: "Invalid signature" }, { status: 400 });

switch (event.event) {
  case "payment.captured": {
    const bookingId = await bookingIdForOrder(payment?.order_id);
    await supabase.from("payments").update({ status: "completed", ... });
    await supabase.from("bookings")
      .update({ status: "confirmed" })
      .eq("id", bookingId)
      .eq("status", "pending");            // only promote PENDING bookings
    await sendBookingConfirmation(bookingId);  // idempotent -- see below
    break;
  }
  ...
}
\end{lstlisting}

Notice \texttt{.eq("status", "pending")} on the update. If a booking was already
cancelled and refunded, a late-arriving webhook must not resurrect it.

\subsection{Idempotency: sending the confirmation email exactly once}

All three paths call \texttt{sendBookingConfirmation()}. Without protection, a
guest could get three identical emails. The solution is an \textbf{atomic claim}:

\begin{lstlisting}[language=TS,caption={lib/emails/bookingConfirmation.ts}]
// Atomic claim: only confirmed bookings that haven't been emailed yet.
const { data: booking } = await admin
  .from("bookings")
  .update({ confirmation_email_sent_at: new Date().toISOString() })
  .eq("id", bookingId)
  .eq("status", "confirmed")
  .is("confirmation_email_sent_at", null)     // <-- the claim
  .select("id, guest_id, ...")
  .maybeSingle();

if (!booking) return;   // already sent, or not confirmed yet
\end{lstlisting}

The \texttt{UPDATE ... WHERE confirmation\_email\_sent\_at IS NULL} is a single
atomic statement. Exactly one caller can win it; every other caller gets zero rows
back and returns early. This is a compare-and-swap, expressed in SQL.

\saythis{Three different code paths can confirm a booking, so all three might try
to email the guest. We make the email idempotent with an atomic claim: a single
UPDATE that stamps a timestamp only if it is still null. Postgres guarantees one
winner. The losers get zero rows and quietly return.}

\subsection{Razorpay Route: splitting the money}

If a hotel manager has linked a Razorpay account, the order carries a
\texttt{transfers} array. Razorpay then automatically settles the manager's share
straight to them, and keeps our commission in our account.

\begin{lstlisting}[language=TS,caption={app/api/payments/razorpay/order/route.ts}]
const { data: payoutAccount } = await supabase.rpc("hotel_payout_account",
  { p_hotel_id: booking.hotel_id });

if (payoutAccount) {
  const basePaise = Math.round(Number(booking.base_price) * 100);
  const managerShare = Math.min(
    Math.round(basePaise * (1 - PLATFORM_COMMISSION_RATE)),  // 15% commission
    amountPaise,                                             // never exceed the charge
  );
  if (managerShare > 0 && managerShare <= amountPaise) {
    transfers = [{ account: payoutAccount, amount: managerShare,
                   currency: "INR", notes: { bookingId } }];
  }
}
\end{lstlisting}

The manager's share is computed on \texttt{base\_price} (the room revenue), not
\texttt{total\_price} --- because GST and the platform fee are not theirs. The
\texttt{Math.min} guard exists because Razorpay rejects an order whose transfers
exceed its amount, which can happen when the hotel only requires an advance
deposit.

\newpage
%==============================================================================
\section{Real-Time Messaging}
\label{sec:messaging}

\subsection{The data model}

\begin{itemize}[leftmargin=*,noitemsep]
  \item \texttt{conversations} --- exactly one per (hotel, guest) pair, enforced by
        \texttt{unique (hotel\_id, guest\_id)}. This is why
        \texttt{getOrCreateConversation()} can safely do
        ``look it up, else insert it''.
  \item \texttt{messages} --- \texttt{body}, \texttt{attachments} (jsonb),
        \texttt{sender\_role} (\texttt{'guest'} or \texttt{'host'}).
\end{itemize}

\subsection{Unread counts are maintained by a trigger}

The browser never increments a counter. Postgres does, on every insert.

\begin{lstlisting}[language=SQL,caption={supabase/hms\_messaging\_migration.sql}]
create or replace function public.on_message_inserted()
returns trigger language plpgsql security definer set search_path = public
as $$
declare v_preview text;
begin
  if new.body is not null and trim(new.body) <> '' then
    v_preview := substring(new.body from 1 for 100);
  else
    v_preview := 'Photo';
  end if;

  update public.conversations
  set last_message_at      = new.created_at,
      last_message_preview = v_preview,
      host_unread  = case when new.sender_role = 'guest' then host_unread  + 1 else host_unread  end,
      guest_unread = case when new.sender_role = 'host'  then guest_unread + 1 else guest_unread end
  where id = new.conversation_id;
  return new;
end; $$;

create trigger on_message_inserted_trigger
  after insert on public.messages
  for each row execute function public.on_message_inserted();
\end{lstlisting}

Because it is a trigger, the counter is correct no matter \emph{how} the message
got inserted --- from the web app, from a script, from the Supabase dashboard.

\subsection{Two kinds of realtime, used for two different things}

Supabase Realtime offers several primitives. We use two, and the distinction is
worth knowing.

\begin{longtable}{@{}p{3.4cm}p{4.4cm}p{7.0cm}@{}}
\toprule
\textbf{Primitive} & \textbf{What it is} & \textbf{Why we chose it here} \\
\midrule
\endhead
\texttt{postgres\_changes} & Streams actual database row changes, read off
Postgres's replication log. & Messages must be durable. If it is in the
database, it is real. \\
\addlinespace
\texttt{broadcast} & An ephemeral pub/sub message that \emph{never touches the
database}. & The ``typing...'' indicator. Writing every keystroke to Postgres
would be absurd. It is transient by nature, so we send it transiently. \\
\bottomrule
\end{longtable}

\begin{lstlisting}[language=TS,caption={components/messaging/useConversationMessages.ts}]
const channel = supabase.channel(`room-${refinedId}`);

channel
  // A. New messages -- durable, from the DB replication stream
  .on("postgres_changes",
    { event: "INSERT", schema: "public", table: "messages",
      filter: `conversation_id=eq.${refinedId}` },
    (payload) => {
      const newMsg = payload.new as Message;
      setMessages((prev) => {
        if (prev.some((m) => m.id === newMsg.id)) return prev;   // de-dupe
        const next = [...prev, newMsg];
        messageCache[activeId] = next;
        return next;
      });
      triggerMarkRead(activeId);
    })
  // B. Typing status -- ephemeral, never stored
  .on("broadcast", { event: "typing" }, (payload) => {
    const { userId, isTyping } = payload.payload;
    if (userId !== currentUserId) setIsOtherPartyTyping(isTyping);
  })
  .subscribe();

return () => { supabase.removeChannel(channel); };   // clean up on unmount
\end{lstlisting}

\begin{notebox}
\textbf{What is a WebSocket?} Normal HTTP is like posting letters: the browser
asks, the server answers, the connection closes. A WebSocket is like leaving the
phone line open. Either side can speak at any time. That is what lets the server
push a message to your screen the instant it is written, with no polling.
\end{notebox}

\begin{warnbox}
\textbf{Realtime respects RLS.} The \texttt{filter} in the subscription is not the
security boundary --- a user could subscribe to another conversation's id. Supabase
Realtime evaluates the \texttt{messages} SELECT policy for the subscribing user
before delivering a row. Since that policy calls
\texttt{can\_access\_conversation()}, an outsider simply receives nothing.
\end{warnbox}

\subsection{Optimistic UI: why sending feels instant}

When you press Enter, the message appears \emph{before} the server has heard about
it.

\begin{lstlisting}[language=TS,caption={components/messaging/useConversationMessages.ts --- sendMessage}]
const tempId = crypto.randomUUID();

// Images render instantly from a local blob URL, before any upload finishes.
const optimisticAttachments = files.map((file) => ({
  url: URL.createObjectURL(file), type: "image",
}));

const optimisticMsg: Message = {
  id: tempId, conversation_id: conversationId, sender_id: currentUserId,
  sender_role: currentUserRole, body, attachments: optimisticAttachments,
  created_at: new Date().toISOString(), read_at: null,
  sending: true,                          // drives the little "clock" icon
};
setMessages((prev) => [...prev, optimisticMsg]);   // <-- appears immediately

try {
  const uploaded = await Promise.all(files.map(uploadToStorage));   // parallel
  const message  = await apiSendMessage(conversationId, body, uploaded, tempId);
  setMessages((prev) => prev.map((m) => (m.id === tempId ? message : m)));
  optimisticAttachments.forEach((a) => URL.revokeObjectURL(a.url));  // free memory
} catch (err) {
  setMessages((prev) => prev.filter((m) => m.id !== tempId));        // roll back
  optimisticAttachments.forEach((a) => URL.revokeObjectURL(a.url));
  throw err;
}
\end{lstlisting}

Two subtleties worth pointing out:

\begin{enumerate}[leftmargin=*,noitemsep]
  \item We pass \texttt{tempId} to the server so the real row gets \emph{the same
        id}. That way when the realtime INSERT event arrives for our own message,
        the \texttt{prev.some(m => m.id === newMsg.id)} de-duplication check
        recognises it and does not draw it twice.
  \item \texttt{URL.revokeObjectURL} is called on both success and failure.
        \texttt{createObjectURL} pins the file in memory until you revoke it ---
        forgetting this is a real memory leak in chat apps.
\end{enumerate}

\subsection{Image compression, before upload}

\begin{lstlisting}[language=TS,caption={lib/compression.ts}]
export async function compressFile(file: File): Promise<File> {
  if (!file.type.startsWith("image/")) return file;   // only images

  const canvas = document.createElement("canvas");
  const MAX_DIMENSION = 1200;                         // plenty for a chat bubble
  // ... scale width/height down, preserving aspect ratio ...
  ctx.drawImage(img, 0, 0, width, height);

  canvas.toBlob(
    (blob) => resolve(new File([blob], name + ".jpg", { type: "image/jpeg" })),
    "image/jpeg",
    0.6,   // quality: heavy compression, still clear
  );
}
\end{lstlisting}

It runs in the browser, so the user's bandwidth is saved, our Storage bill is
saved, and the recipient's load is instant. A 12~MB phone photo becomes roughly
150~KB.

\subsection{The privacy problem, and \texttt{shares\_conversation\_with}}

The chat panel shows Call / WhatsApp / Email buttons for the host. Those need the
host's phone number, which lives in \texttt{profiles} --- a table whose RLS says
``you may only read your own row''.

The naive fix (``let guests read manager profiles'') would turn the platform into
a phone directory. Instead we wrote a function that answers one narrow question.

\begin{lstlisting}[language=SQL,caption={supabase/hms\_messaging\_names\_migration.sql}]
create or replace function public.shares_conversation_with(p_profile_id uuid)
returns boolean language sql stable
security definer set search_path = public
as $$
  select exists (
    select 1 from public.conversations c
    where c.guest_id = p_profile_id
      and ( c.guest_id = auth.uid()
         or public.can_manage_hotel(c.hotel_id, 'reply_messages')
         or public.get_my_role() = 'admin' )
  );
$$;

create policy "profiles: read conversation guests"
  on public.profiles for select
  using (public.shares_conversation_with(id));
\end{lstlisting}

So contact details become visible \emph{only} once a real conversation exists
between those two parties, and only to a party with permission to reply in it.

For the reverse direction (public pages showing reviewer names), we do not expose
\texttt{profiles} at all. We expose a \emph{view}:

\begin{lstlisting}[language=SQL,caption={supabase/listing\_public\_migration.sql}]
-- Exposes ONLY id, full_name, created_at. Never phone, dob, email, or location.
CREATE OR REPLACE VIEW public.public_profiles AS
  SELECT id, full_name, created_at FROM public.profiles;

GRANT SELECT ON public.public_profiles TO anon, authenticated;
\end{lstlisting}

\saythis{A guest needs the host's phone number to call them --- but only if they are
actually talking to that host. So instead of loosening the profiles policy, we
wrote a \texttt{SECURITY DEFINER} function that checks for a shared conversation,
and made \emph{that} the policy. For public pages we never touch \texttt{profiles};
we expose a view containing only name and join date.}

\subsection{The contact buttons themselves}

No library. Just URL schemes every device understands.

\begin{lstlisting}[language=HTML]
<a href="tel:+917736184696">Call</a>
<a href="https://wa.me/917736184696">WhatsApp</a>   <!-- digits only, no + or spaces -->
<a href="mailto:host@example.com">Email</a>
\end{lstlisting}

\newpage
%==============================================================================
\section{QR Code Check-In}

\subsection{What the QR code actually contains}

Not a secret. Not a token. A plain URL:

\begin{lstlisting}[language=TS,caption={components/SuccessClient.tsx / app/bookings/[id]/page.tsx}]
import QRCode from "qrcode";

const qrCodeUrl = await QRCode.toDataURL(
  `${origin}/bookings/${booking.id}/check-in`
);
// -> "data:image/png;base64,iVBORw0KGgo..."  (an image, inline, no file needed)
\end{lstlisting}

\begin{notebox}
\textbf{Why is it safe to put the booking id in a QR code?} Because knowing an id
grants nothing. To \emph{use} it, the scanner must call
\texttt{check\_in\_booking()}, which asks Postgres
\texttt{can\_operate\_hotel(booking.hotel\_id)}. Only that hotel's manager or
assigned staff pass. A stranger who photographs your QR code can look at the
booking page (RLS blocks them) and check you in (the RPC blocks them). Nothing.

Note also that the QR encodes a \emph{URL}, so a guest who scans their own code
with a normal camera app lands on their check-in page. Nice touch, and it means
the code degrades gracefully.
\end{notebox}

Note that \texttt{toDataURL} returns a base64 data URI, which we drop straight
into an \texttt{<img src>}. No file is written, nothing is uploaded, nothing is
stored. The QR is regenerated on every page render from the booking id.

\subsection{Scanning it (front desk)}

\begin{lstlisting}[language=TS,caption={app/manager/manage/[hotelId]/FrontDesk.tsx}]
import { Html5Qrcode } from "html5-qrcode";

useEffect(() => {
  const html5Qrcode = new Html5Qrcode("qr-reader");
  html5Qrcode.start(
    { facingMode: "environment" },        // rear camera on a phone
    { fps: 10, qrbox: (w, h) => { const s = Math.min(w, h) * 0.7;
                                  return { width: s, height: s }; } },
    (decodedText) => onScan(decodedText),
    () => { /* ignore scan failures -- normal when the QR is out of frame */ },
  );
  return () => { if (html5Qrcode.isScanning) html5Qrcode.stop(); };  // release camera
}, [onScan, onError]);
\end{lstlisting}

The cleanup function is not optional. If you unmount without calling
\texttt{stop()}, the camera light stays on.

\subsection{From scanned text to a check-in}

\begin{lstlisting}[language=TS,caption={app/manager/manage/[hotelId]/FrontDesk.tsx}]
onScan={async (decodedText) => {
  // The QR holds a URL; pull the UUID out of it. Fall back to raw text so a
  // plain booking id also scans.
  const match = decodedText.match(/bookings\/([a-f0-9-]{36})/i);
  const bookingId = match ? match[1] : decodedText.trim();

  let matched = bookings.find((b) => b.id.toLowerCase() === bookingId.toLowerCase());

  // The booking may be newer than our cached list -- re-fetch before giving up.
  if (!matched) {
    const res = await fetch(`/api/manager/manage/${hotel.id}`);
    ...
  }
}}
\end{lstlisting}

Then the server action:

\begin{lstlisting}[language=SQL,caption={supabase/hms\_management\_migration.sql}]
create or replace function public.check_in_booking(p_booking_id uuid)
returns void language plpgsql security definer set search_path = public
as $$
declare v_b public.bookings;
begin
  select * into v_b from public.bookings where id = p_booking_id;
  if not found then raise exception 'Booking not found'; end if;
  if not public.can_operate_hotel(v_b.hotel_id) then raise exception 'Not allowed'; end if;
  if v_b.status <> 'confirmed' then
    raise exception 'Only confirmed bookings can be checked in'; end if;
  update public.bookings set status = 'checked_in' where id = p_booking_id;
end; $$;
\end{lstlisting}

Three guards, in order: does it exist, may you touch this hotel, is it in the
right state. A pending (unpaid) booking cannot be checked in. A cancelled one
cannot. And it cannot be checked in twice, because after the first time its
status is no longer \texttt{'confirmed'}.

\saythis{The QR is just a URL containing the booking id. Security is not in the
code --- it is in \texttt{check\_in\_booking}, which verifies the scanner is that
hotel's manager or staff via \texttt{can\_operate\_hotel}, and that the booking is
in \texttt{confirmed} state. Scan it twice and the second one is rejected, because
the status has moved on.}

\newpage
%==============================================================================
\section{Maps}

\subsection{What we use, and what we deliberately did not}

\begin{longtable}{@{}p{3.2cm}p{11.6cm}@{}}
\toprule
\textbf{Option} & \textbf{Verdict} \\
\midrule
\endhead
Google Maps & \textbf{Rejected.} Requires an API key tied to a credit card. The
key must be shipped to the browser, where anyone can steal it and run up your
bill. Free tier ends. \\
\addlinespace
Mapbox & \textbf{Rejected.} Same key-in-the-browser problem, plus a monthly
map-load quota. \\
\addlinespace
Static PNG tiles & \textbf{Rejected.} Not interactive. \\
\addlinespace
\textbf{MapLibre GL + OpenFreeMap} & \textbf{Chosen.} MapLibre is the open-source
fork of Mapbox GL (from before it went proprietary). OpenFreeMap serves
OpenStreetMap vector tiles with \textbf{no API key and no rate limit}. Renders
with WebGL, so panning and zooming are buttery. \\
\addlinespace
\textbf{Leaflet + CartoDB Voyager} & \textbf{Chosen, as the fallback.} If WebGL is
unavailable or the vector style fails to load, we transparently swap to Leaflet
with raster tiles. The map never disappears. \\
\bottomrule
\end{longtable}

\begin{lstlisting}[language=TS,caption={components/LocationMap.tsx}]
// Rich, free, no-key OpenStreetMap vector tiles (ODbL). "positron" is the clean,
// muted, modern style -- keeps street/place labels but drops the loud OSM colors.
const OPENFREEMAP_STYLE = "https://tiles.openfreemap.org/styles/positron";

// CartoDB Voyager -- the richest of our free raster basemaps. Used as the
// fallback engine (Leaflet) when the vector map fails to load.
const VOYAGER_TILES = "https://{s}.basemaps.cartocdn.com/rastertiles/voyager/{z}/{x}/{y}{r}.png";

/**
 * One interactive map surface. Tries MapLibre + OpenFreeMap (rich vector); if
 * that fails to load, transparently falls back to Leaflet + Voyager so the map
 * never disappears. Both engines share the same pin, popup and zoom controls.
 */
\end{lstlisting}

\subsection{Vector tiles vs raster tiles}

\textbf{Raster} tiles are pre-rendered PNG images, 256$\times$256 pixels. Zoom in
and they go blurry until the next zoom level downloads.

\textbf{Vector} tiles ship the \emph{geometry} --- ``this road is a line through
these coordinates'' --- and the browser draws it with WebGL. So labels stay sharp
at any zoom, the map can rotate and tilt, and you can restyle it (that is how we
get the muted ``positron'' palette that matches our cream-and-green theme without
touching a server).

\subsection{Where maps appear}

\begin{itemize}[leftmargin=*,noitemsep]
  \item \file{app/hotels/[id]/components/LocationSection.tsx} --- one pin, the
        hotel, with a fullscreen mode.
  \item \file{app/hotels/ExploreMap.tsx} --- every approved hotel as a clickable
        pin, split-screen with the results grid.
  \item \file{app/manager/create-hotel/components/LeafletMap.tsx} --- the
        \emph{picker}. The manager clicks to drop a pin; we store
        \texttt{latitude} and \texttt{longitude} on the hotel row. Leaflet is used
        here because click-to-place needs nothing fancy, and it is smaller.
\end{itemize}

\saythis{We use MapLibre with OpenFreeMap vector tiles --- both fully open source,
no API key, no billing. Google Maps and Mapbox both require a key that has to be
shipped to the browser, where anyone can lift it. And we kept Leaflet with raster
tiles as an automatic fallback, so if WebGL is unavailable the map still renders.}

\newpage
%==============================================================================
\section{Email Notifications}

\subsection{Why Brevo, and why HTTP instead of SMTP}

\begin{lstlisting}[language=TS,caption={lib/email.ts --- the header comment is the answer}]
// Transactional email via Brevo's HTTP API. HTTP (not SMTP) so it works
// reliably on serverless. Server-only: never import this into client code,
// it reads BREVO_API_KEY. All sends are best-effort and never throw -- a failed
// email must never break the action that triggered it.
\end{lstlisting}

Three reasons, all in that comment:

\begin{enumerate}[leftmargin=*,noitemsep]
  \item \textbf{HTTP, not SMTP.} Serverless functions on Vercel are short-lived
        and often block outbound SMTP ports. An HTTPS \texttt{fetch} always works.
  \item \textbf{Server-only.} It reads \texttt{BREVO\_API\_KEY}. Importing it into
        a Client Component would leak the key.
  \item \textbf{Never throws.} \texttt{sendEmail} returns \texttt{false} on
        failure. A down mail provider must never fail a paid booking.
\end{enumerate}

And why Brevo over SendGrid/Mailgun/Resend: Brevo's free tier does not require you
to own and verify a custom domain, which matters for a student project with no
domain to verify.

\begin{lstlisting}[language=TS,caption={lib/email.ts}]
export async function sendEmail(input: SendEmailInput): Promise<boolean> {
  const apiKey = process.env.BREVO_API_KEY;
  if (!apiKey || !from) {
    console.error("[email] BREVO_API_KEY or EMAIL_FROM missing -- skipping send");
    return false;
  }
  try {
    const res = await fetch("https://api.brevo.com/v3/smtp/email", {
      method: "POST",
      headers: { "api-key": apiKey, "content-type": "application/json" },
      body: JSON.stringify({
        sender: { email: from, name: fromName },
        to: [{ email: input.to, name: input.toName }],
        subject: input.subject,
        htmlContent: input.html,
      }),
    });
    if (!res.ok) { console.error("[email] Brevo send failed", res.status); return false; }
    return true;
  } catch (err) { console.error("[email] send error", err); return false; }
}
\end{lstlisting}

\subsection{The emails we send}

One file each, under \file{lib/emails/}:
\texttt{bookingConfirmation}, \texttt{bookingCancellation},
\texttt{checkInConfirmation}, \texttt{checkOutConfirmation},
\texttt{refundProcessed}, \texttt{staffInvite}, \texttt{managerVerification},
\texttt{listingDecision}.

\subsection{Why email HTML is written with tables and inline styles}

If you open \file{lib/email.ts} you will find \texttt{<table role="presentation">}
everywhere and every style written inline. This looks like 2003 HTML, and it is
deliberate: \textbf{email clients strip \texttt{<style>} blocks}, and Outlook renders
with Microsoft Word's layout engine, which does not understand flexbox or grid.
Tables and inline styles are the only things that render identically in Gmail,
Outlook, and Apple Mail. The comment in the source says exactly this:
\emph{``Inline styles only (email clients strip \texttt{<style>}/CSS).''}

The QR code in the confirmation email is fetched from an external QR image service
rather than embedded, because many clients block base64 images.

\newpage
%==============================================================================
\section{Managers, Staff and the Admin Panel}

\subsection{Staff are real accounts, not sub-users}

A staff member has \texttt{role = 'staff'} in \texttt{profiles} and a row in
\texttt{hotel\_staff} linking them to \emph{one hotel} with a
\texttt{permissions text[]} array (e.g.
\texttt{\{offline\_booking, view\_occupancy, reply\_messages\}}).

Every permission check in the database funnels through one of two functions:

\begin{itemize}[leftmargin=*,noitemsep]
  \item \texttt{can\_operate\_hotel(hotel)} --- manager \emph{or any} staff.
        Used for check-in/check-out.
  \item \texttt{can\_manage\_hotel(hotel, perm)} --- manager \emph{or} staff who
        hold that specific permission. Used for messaging and room management.
\end{itemize}

\subsection{The invite flow}

\begin{enumerate}[leftmargin=*,noitemsep]
  \item Manager enters an email and ticks permissions.
  \item A row is created in \texttt{staff\_invites} with a random UUID
        \texttt{token} and \texttt{expires\_at = now() + 14 days}.
  \item Brevo sends the invitation link.
  \item The invitee signs up (or logs in) and lands on \file{app/staff/accept/}.
  \item \texttt{accept\_my\_staff\_invites()} matches pending invites by the
        \emph{logged-in user's email}, creates the \texttt{hotel\_staff} rows, and
        flips their profile role to \texttt{staff}.
\end{enumerate}

Matching on the authenticated email (not just the token) means a stolen link
cannot be redeemed by somebody else's account.

\subsection{Manager verification}

A manager cannot list a hotel until an admin approves them. \file{proxy.ts}
enforces this: any \texttt{/manager/*} request from an unapproved manager is
redirected to \texttt{/manager/waiting}. Statuses are \texttt{pending},
\texttt{approved}, \texttt{rejected}, and \texttt{more\_info} --- the last lets an
admin ask for a better document without a hard rejection.

\subsection{Every admin action is an audited RPC}

\begin{lstlisting}[language=SQL,caption={supabase/hms\_admin\_migration.sql}]
create or replace function public.admin_review_hotel(
  p_id uuid, p_decision text, p_reason text default null)
returns void language plpgsql security definer set search_path = public
as $$
declare v_uid uuid := auth.uid(); v_h public.hotels;
begin
  if public.get_my_role() <> 'admin' then raise exception 'Not allowed'; end if;
  if p_decision not in ('approved', 'rejected') then raise exception 'Invalid decision'; end if;

  update public.hotels
    set status = p_decision,
        rejection_reason = case when p_decision = 'rejected' then p_reason else null end,
        reviewed_at = now(), reviewed_by = v_uid,
        published_at = case when p_decision = 'approved' then now() else published_at end
    where id = p_id;

  insert into public.admin_audit_log (admin_id, action, target_type, target_id, details)
    values (v_uid, 'hotel_' || p_decision, 'hotel', p_id,
            jsonb_build_object('reason', p_reason, 'name', v_h.name));
end; $$;
\end{lstlisting}

The pattern never varies: \emph{check the role, validate the input, do the work,
write the audit row} --- all in one transaction, so an action can never be
performed without being logged.

\saythis{Admin powers do not live in the API layer. They live in
\texttt{SECURITY DEFINER} functions that begin by asserting
\texttt{get\_my\_role() = 'admin'} and end by writing an audit-log row --- in the
same transaction. So an admin action that is not logged is not possible.}

\newpage
%==============================================================================
\section{Every Package, and Why It Is There}

\subsection{Runtime dependencies}

\begin{longtable}{@{}p{3.3cm}p{4.7cm}p{6.8cm}@{}}
\toprule
\textbf{Package} & \textbf{What it does} & \textbf{Why this one} \\
\midrule
\endhead
\texttt{next} 16.2 & The framework: routing, SSR, API routes, caching. &
One codebase for frontend and backend. Server Components remove a whole class of
API endpoints. \\
\addlinespace
\texttt{react} 19.2 & The UI library. & Required by Next.js. v19 gives us Server
Components and Actions. \\
\addlinespace
\texttt{@supabase/supabase-js} & The database/auth/storage/realtime client. &
The official SDK. \\
\addlinespace
\texttt{@supabase/ssr} & Cookie-aware Supabase clients for server rendering. &
Without it, the server cannot read the user's session, so RLS would see every
server-side query as anonymous. \\
\addlinespace
\texttt{razorpay} & Server-side SDK for creating orders and fetching payments. &
Only used server-side --- it holds the secret. \\
\addlinespace
\texttt{swr} & Client data fetching with stale-while-revalidate + preloading. &
Tiny (about 4~KB), made by the Next.js team, and its \texttt{preload()} gives us
hover-prefetching on the dashboards. Chosen over TanStack Query because we needed
caching, not a full query engine. \\
\addlinespace
\texttt{maplibre-gl} & WebGL vector map renderer. & Open-source, no API key. \\
\addlinespace
\texttt{leaflet} & Lightweight raster map. & Fallback engine + the location picker
in the create-hotel wizard. \\
\addlinespace
\texttt{qrcode} & Generates the check-in QR as a base64 data URI. & Runs on the
server; no image file to store. \\
\addlinespace
\texttt{html5-qrcode} & Reads QR codes from a device camera. & Browser-only;
handles camera permissions and decoding. \\
\addlinespace
\texttt{isomorphic-dompurify} & Sanitises HTML. & The rich-text hotel description
is manager-supplied HTML. Rendering it raw would be a stored-XSS hole. This
version runs in \emph{both} browser and Node, so we sanitise on write
(\file{lib/hotelDraft.ts}) \emph{and} on render
(\file{OverviewSection.tsx}). \\
\addlinespace
\texttt{motion} & Animation (the library formerly called Framer Motion). &
Used in 37 files for page transitions, modals, and the chat. Declarative:
\texttt{animate=\{...\}} instead of manual CSS keyframes. \\
\addlinespace
\texttt{gsap} & Timeline animation. & Used once, for the landing-page hero
sequence, where we needed precisely choreographed timing. \\
\addlinespace
\texttt{lucide-react} & SVG icon set. & Tree-shakeable; each icon is a React
component, so we ship only the icons we use. \\
\addlinespace
\texttt{@animateicons/react} & Animated icons. & Used in the navbar. \\
\addlinespace
\texttt{canvas-confetti} & The confetti burst on booking success. & Pure joy.
About 3~KB. \\
\bottomrule
\end{longtable}

\subsection{Development dependencies}

\begin{longtable}{@{}p{3.3cm}p{11.5cm}@{}}
\toprule
\textbf{Package} & \textbf{Purpose} \\
\midrule
\endhead
\texttt{typescript} & Static types. Catches ``you forgot this field'' before
runtime. Everything shared lives in \file{lib/types.ts}. \\
\addlinespace
\texttt{tailwindcss} v4 + \texttt{@tailwindcss/postcss} & Utility-first CSS.
v4 configures itself in \file{app/globals.css} --- there is no
\texttt{tailwind.config.js} any more. \\
\addlinespace
\texttt{eslint}, \texttt{eslint-config-next} & Linting with the Next.js rule set
(catches, e.g., using a hook in a Server Component). \\
\addlinespace
\texttt{pg} & A raw Postgres driver, used only by the local seed script
(\file{supabase/seed\_hotels.ts}). Not shipped. \\
\addlinespace
\texttt{@types/*} & Type definitions for the JS-only libraries above. \\
\bottomrule
\end{longtable}

\subsection{Things we chose \emph{not} to use}

\begin{longtable}{@{}p{3.3cm}p{11.5cm}@{}}
\toprule
\textbf{Not used} & \textbf{Why not} \\
\midrule
\endhead
Redux / Zustand & There is very little global client state. Server Components hold
the data; \texttt{AuthProvider} and \texttt{CurrencyProvider} are two small React
Contexts. A store would have been ceremony without benefit. \\
\addlinespace
Prisma / Drizzle & An ORM would sit \emph{above} the database and would tempt us
into writing authorisation in TypeScript. We wanted RLS and PL/pgSQL to be the
source of truth, so we talk to Postgres through PostgREST and RPC. \\
\addlinespace
Express / Nest & Next.js Route Handlers already are the backend. A second server
would need its own deploy, its own auth, its own CORS. \\
\addlinespace
Socket.io & Would require a long-lived server process, which Vercel's serverless
functions are not. Supabase Realtime is a managed WebSocket layer that already
knows our RLS policies. \\
\addlinespace
Stripe & Weak UPI support in India. Razorpay is the local standard. \\
\bottomrule
\end{longtable}

\newpage
%==============================================================================
\section{Security Summary}

If you are asked ``how secure is this?'', walk down this list.

\begin{enumerate}[leftmargin=*]
  \item \textbf{Row Level Security on every table.} The database refuses to return
        rows you may not see. Not the API --- the database.
  \item \textbf{The service-role key is used in exactly one file}
        (\file{lib/supabase/admin.ts}) and imported by exactly one caller: the
        Razorpay webhook, where no user session exists.
  \item \textbf{Every \texttt{SECURITY DEFINER} function pins
        \texttt{search\_path = public}}, closing the classic privilege-escalation
        path.
  \item \textbf{Prices are never accepted from the client.} \texttt{book\_room}
        reads \texttt{rooms.price} itself.
  \item \textbf{Payment signatures are HMAC-verified with
        \texttt{crypto.timingSafeEqual}}, both for the checkout callback and the
        webhook (the webhook signs the \emph{raw body}, so we read
        \texttt{request.text()} before parsing).
  \item \textbf{Manager-supplied HTML is sanitised with DOMPurify}, on write and
        on render.
  \item \textbf{Overbooking is impossible} by CHECK constraint, independent of
        application logic.
  \item \textbf{Emails and payment confirmations are idempotent}, so a retried
        webhook cannot double-charge or double-send.
  \item \textbf{Personal data is never exposed by default.} Public pages read
        \texttt{public\_profiles} (id, name, join date). Phone numbers unlock only
        via \texttt{shares\_conversation\_with}.
  \item \textbf{Every admin action is written to an append-only audit log}, in the
        same transaction as the action.
\end{enumerate}

\newpage
%==============================================================================
\section{Frequently Asked Questions}

Grouped by the part of the system they attack. Read the question, try to answer
out loud, \emph{then} read the answer.

\subsection{Architecture and Next.js}

\Q{Why Next.js and not plain React with an Express backend?}
\A{Because we would then be maintaining two applications, two deploys, and a CORS
configuration. With Next.js the same project renders the HTML \emph{and} serves
the API routes. It also gives us Server Components, which let a page read the
database directly and ship zero JavaScript for that work --- with plain React,
every page would need a loading spinner and an API endpoint.}

\Q{What is the difference between a Server Component and a Client Component?}
\A{A Server Component runs on the server only. Its JavaScript is never sent to the
browser, and it can \texttt{await} a database query directly. A Client Component
(marked \texttt{"use client"}) is shipped to the browser and can use state,
effects and browser APIs. Our rule: fetch in the Server Component, pass the data
as props to a small Client Component for interactivity. That is why the codebase
is full of \texttt{page.tsx} plus \texttt{SomethingClient.tsx} pairs.}

\Q{How does routing work? Where is the route config?}
\A{There is none. The folder structure \emph{is} the routing table.
\file{app/hotels/[id]/page.tsx} serves \texttt{/hotels/<id>}. A folder in square
brackets is a dynamic segment. A \file{route.ts} makes the folder an API endpoint
rather than a page.}

\Q{Why is your middleware file called \file{proxy.ts}?}
\A{Next.js 16 renamed the convention from \file{middleware.ts} to \file{proxy.ts}.
It does the same job --- runs before matching requests. Our \texttt{matcher} only
lists \texttt{/dashboard}, \texttt{/manager} and \texttt{/admin}, so public pages
skip it entirely and pay no auth cost.}

\Q{If middleware protects your routes, what happens if I call the API directly?}
\A{Nothing bad. \texttt{/api} is deliberately \emph{not} in the matcher. Every API
route authenticates independently --- admin routes call \texttt{getAdminContext()},
and every query is subject to RLS regardless. Middleware exists so users get a
friendly redirect, not as the security boundary.}

\Q{What is a Server Action, and is it safe?}
\A{A function marked \texttt{"use server"} that a Client Component can import and
call like a normal function; Next.js turns the call into a POST behind the scenes.
It is \emph{not} inherently safe --- it is a public endpoint, and anyone can call
it with any arguments. That is why every action in \file{app/messages/actions.ts}
and friends re-verifies the user with \texttt{auth.getUser()} and relies on RLS.}

\Q{Explain your caching strategy.}
\A{Three layers. (1) \texttt{unstable\_cache} caches the expensive approved-hotels
query result under the tag \texttt{"hotels"} with a one-hour fallback. (2) When an
admin approves or a manager edits a hotel, we call \texttt{revalidateTag("hotels")}
and the cache is dropped immediately. (3) The hotel detail page also sets
\texttt{export const revalidate = 60}, so its rendered HTML is CDN-cached for a
minute. Fast by default, correct the moment data changes.}

\Q{Why is there a \file{lib/supabase/public.ts} when \file{server.ts} exists?}
\A{Because \texttt{unstable\_cache} forbids reading request-scoped data such as
cookies --- the cached value must not depend on who asked. \file{server.ts} reads
cookies. So for cached queries we use a stateless, cookie-free anon client. RLS
still applies; anonymous can only read approved hotels.}

\Q{What is SWR and why did you use it over TanStack Query?}
\A{SWR shows the cached (stale) data instantly, then revalidates in the background.
We picked it because it is about 4~KB, comes from the Next.js team, and its
\texttt{preload()} lets us prefetch a dashboard's data when the user merely
\emph{hovers} its sidebar link --- see \texttt{managerApiForRoute()} in
\file{lib/swr.ts}. We needed a cache, not a full query engine.}

\subsection{Supabase and the database}

\Q{What is Row Level Security, in one sentence?}
\A{A rule attached to a table that Postgres silently appends to every query, so a
user physically cannot read or write rows the rule excludes --- even if they hold a
valid API key and bypass our application entirely.}

\Q{Isn't it dangerous to put a database key in the browser?}
\A{No --- that is the design. The anon key gets you \emph{to} the database; RLS
decides what you get \emph{out of} it. The dangerous key is the service-role key,
which bypasses RLS. That one lives in \file{lib/supabase/admin.ts}, is never
prefixed \texttt{NEXT\_PUBLIC\_}, and is used only by the Razorpay webhook.}

\Q{Why do you have four Supabase clients?}
\A{Browser (manages cookies), server (reads this request's cookies so RLS knows
who you are), public (cookie-free, because caching forbids request-scoped data),
and admin (service role, bypasses RLS, webhook only). Four constraints, four
clients.}

\Q{You have a lot of logic in SQL. Why not in TypeScript?}
\A{Four reasons. Atomicity: \texttt{book\_room} reserves inventory, creates the
booking and creates the payment in one transaction that rolls back entirely on
failure. Concurrency: only the database can take a row lock. Trust: the price is
computed from the database's own numbers, so a tampered client cannot set it. And
single-source: the web app and the webhook both confirm payments by calling the
same function.}

\Q{What is \texttt{SECURITY DEFINER} and why is it everywhere?}
\A{It makes a function run with the privileges of its creator rather than its
caller, so RLS does not apply inside it. We need it for two things: to let a
policy read a table without triggering that table's own policy (avoiding infinite
recursion), and to let a guest trigger a controlled write they could not perform
directly. Every one of ours pins \texttt{search\_path = public} so an attacker
cannot swap in a malicious schema.}

\Q{You mentioned RLS recursion. Show me the actual problem.}
\A{The \texttt{hotels} policy wants to know ``is the caller staff on this hotel?'',
which reads \texttt{hotel\_staff}. The \texttt{hotel\_staff} policy wants to know
``does the caller own this hotel?'', which reads \texttt{hotels}. Each policy
triggers the other, forever. Postgres throws \texttt{infinite recursion detected
in policy}. We broke the cycle with single-table \texttt{SECURITY DEFINER} helpers
--- \texttt{owns\_hotel()} and \texttt{is\_hotel\_staff()} --- which read one table
each with RLS bypassed. That is what \file{hms\_rls\_recursion\_fix.sql} is.}

\Q{How does a user get a profile row?}
\A{A trigger, \texttt{handle\_new\_user}, fires \texttt{after insert on
auth.users}. It reads the sign-up metadata we passed to
\texttt{supabase.auth.signUp()} and inserts the profile --- and, if the role is
manager, a pending \texttt{manager\_verifications} row too. Because it is a
trigger, a user without a profile is structurally impossible.}

\Q{How do you make hotel search fast?}
\A{Searching is \texttt{ILIKE '\%goa\%'}, and a normal B-tree index cannot serve a
leading wildcard --- Postgres would scan every row. So we enabled the
\texttt{pg\_trgm} extension and created GIN trigram indexes on name, location,
city and state. We also created a \emph{partial} index on
\texttt{(created\_at desc) where status = 'approved'}, which satisfies both the
filter and the sort of the catalogue query with no sort step.}

\Q{Why is there a \texttt{public\_profiles} view?}
\A{Because review author names and the host card are public, but phone numbers and
dates of birth are not. The view selects only \texttt{id}, \texttt{full\_name} and
\texttt{created\_at}, and \emph{that} is what we grant to \texttt{anon}. There is
no way to widen it accidentally.}

\subsection{Booking and concurrency}

\Q{Two users book the last room at the same instant. What happens?}
\A{This is the question I most wanted to be asked. Inside \texttt{book\_room} we do
\texttt{SELECT total\_count - booked\_count ... FOR UPDATE}, which takes a
row-level lock on that night's inventory row. The second transaction blocks on that
lock, and when it resumes it reads the \emph{updated} count, sees zero available,
and raises an exception --- which rolls back the whole function, including any
nights it had already reserved. And as a final backstop the table has
\texttt{check (booked\_count <= total\_count)}, so even a bug in our logic could
not overbook the hotel.}

\Q{Why is availability \texttt{total\_units - max(booked\_count)} and not the sum
or the average?}
\A{Because a guest needs the room on \emph{every} night of their stay. If the room
is full on the middle night, the whole range is unbookable. Taking the maximum
booked count across the range gives exactly the number of units free for the
entire stay.}

\Q{Why \texttt{i.date < p\_check\_out} rather than \texttt{<=}?}
\A{You do not occupy a room on the night of the day you check out. A stay from the
14th to the 16th consumes the nights of the 14th and 15th --- two nights, two
inventory rows.}

\Q{Do you create an inventory row for every room for every future date?}
\A{No, that would be tens of thousands of rows per hotel. They are created lazily:
\texttt{insert ... on conflict (room\_id, date) do nothing} runs just before we
lock the row, so the row exists exactly when it is first needed.}

\Q{You calculate the price in TypeScript \emph{and} in SQL. Isn't that a bug
waiting to happen?}
\A{The TypeScript one is a preview so the breakdown updates instantly as you pick
dates --- and the source comment says so: ``Mirror of the server-side price math in
\texttt{book\_room} (kept in sync for preview).'' Its number is never sent to the
server. \texttt{POST /api/bookings} accepts a room id and dates, nothing else; the
database computes the real price from \texttt{rooms.price} and
\texttt{hotels.gst\_percent}. If they drifted, the user would see a wrong preview,
never a wrong charge.}

\Q{What is your refund policy and where does it live?}
\A{100\% at 48 hours or more before check-in, 50\% between 24 and 48 hours,
nothing under 24 hours. It is computed in \texttt{cancel\_booking}, which also
releases the inventory night by night using \texttt{greatest(0, booked\_count -
num\_rooms)} so the count can never go negative.}

\Q{Why does \texttt{cancel\_booking} appear in two SQL files?}
\A{Migrations are applied in order and each uses \texttt{create or replace}. The
first version (in \file{booking\_schema.sql}) allowed the guest and admins. The
later one (in \file{hms\_management\_migration.sql}) also allows the hotel's
manager and staff, so the front desk can cancel for a guest who phones the hotel.
The later definition wins.}

\subsection{Payments}

\Q{Walk me through a payment.}
\A{The browser POSTs a booking id to \file{/api/payments/razorpay/order}. The
server reads the booking (RLS proves ownership), converts the total to paise, calls
\texttt{razorpay.orders.create}, stores the returned order id on the payment row,
and returns the order id plus the \emph{publishable} key. The browser opens
Razorpay's checkout modal --- card details go straight to Razorpay, never to us. On
success Razorpay hands the browser a payment id and a signature, which we POST to
\file{/api/payments/razorpay/verify}. There we recompute
\texttt{HMAC\_SHA256(order\_id + "|" + payment\_id)} with our secret and compare it
in constant time. If it matches, we call \texttt{confirm\_razorpay\_payment} and
the booking becomes confirmed.}

\Q{Why verify on the server? The browser already knows the payment succeeded.}
\A{Because the browser is not trustworthy. Anyone can open dev tools and POST
\texttt{\{bookingId, paymentId: "anything"\}} to our verify endpoint. The signature
is the proof: it is an HMAC keyed with a secret only Razorpay and our server hold.
Without the secret you cannot forge it.}

\Q{What is \texttt{timingSafeEqual} and why not just \texttt{===}?}
\A{\texttt{===} short-circuits at the first differing character, so the time it
takes leaks how many leading characters were correct. With enough precise
measurements an attacker can recover the signature one byte at a time. That is a
timing attack. \texttt{crypto.timingSafeEqual} always compares every byte.}

\Q{What if the user's browser dies right after paying?}
\A{Three independent paths confirm a booking, and we need all three. The happy
path is the checkout callback. The webhook is Razorpay's server calling ours
directly --- it works even with the browser closed, and because there is no user
session it uses the service-role client. And the reconcile path runs when the user
re-opens the booking: it asks Razorpay \texttt{orders.fetchPayments(orderId)} and
confirms if a captured payment exists.}

\Q{If three paths can confirm the booking, does the guest get three emails?}
\A{No. The email function performs an atomic claim: a single
\texttt{UPDATE bookings SET confirmation\_email\_sent\_at = now() WHERE id = \$1
AND status = 'confirmed' AND confirmation\_email\_sent\_at IS NULL} that returns
the row. Postgres guarantees exactly one caller wins; the others get zero rows and
return immediately. It is a compare-and-swap written in SQL.}

\Q{Why does the webhook read \texttt{request.text()} instead of
\texttt{request.json()}?}
\A{Because the signature is computed over the \emph{exact bytes} Razorpay sent.
Parsing to JSON and re-serialising would change whitespace and key order, and the
HMAC would no longer match. You must hash the raw body.}

\Q{The webhook updates a booking to confirmed. What stops it resurrecting a
cancelled booking?}
\A{The update is
\texttt{.eq("id", bookingId).eq("status", "pending")}. Only a booking still awaiting
payment can be promoted. A late webhook for a booking that has since been cancelled
matches zero rows.}

\Q{Why Razorpay and not Stripe?}
\A{The platform is India-first. Razorpay has first-class UPI, net banking and
Indian card support, prices in INR natively, and has a proper test mode. Stripe's
UPI support in India is weak. Razorpay also gave us \emph{Route}, which lets a
single payment be split so the hotel's share settles directly to their account
while we keep our 15\% commission.}

\Q{How does the commission split work?}
\A{If the manager has linked a Razorpay account, the order carries a
\texttt{transfers} array with \texttt{base\_price * (1 - 0.15)} in paise. Razorpay
settles that to them and keeps the rest with us. Note it is computed on
\texttt{base\_price}, the room revenue --- not on \texttt{total\_price}, because GST
and the platform fee are not the manager's money. And we clamp it with
\texttt{Math.min(managerShare, amountPaise)} because Razorpay rejects an order
whose transfers exceed its own amount, which can happen when a hotel only takes an
advance deposit.}

\Q{Are you handling real money?}
\A{We are in Razorpay \emph{test mode}, which simulates the full flow --- including
webhooks and signature verification --- with test card numbers, so no real money
moves. Switching to production is a matter of changing the API keys; not a line of
our code changes.}

\subsection{Messaging}

\Q{How is the chat real-time? Are you polling?}
\A{No polling. The browser opens a WebSocket to Supabase Realtime and subscribes to
\texttt{postgres\_changes} for \texttt{INSERT} on the \texttt{messages} table,
filtered to one conversation. Supabase reads Postgres's replication log, so the
moment a row is committed it is pushed to every subscriber.}

\Q{Someone could subscribe to another conversation's id. Isn't that a data leak?}
\A{No. The filter is a convenience, not a security control. Supabase Realtime
evaluates the \texttt{messages} SELECT policy for the subscribing user before
delivering any row, and that policy calls \texttt{can\_access\_conversation()}. An
outsider's subscription simply receives nothing.}

\Q{Why is the typing indicator not stored in the database?}
\A{Because it is worthless one second later. We send it as a Realtime
\texttt{broadcast}, which is ephemeral pub/sub that never touches Postgres.
Writing every keystroke to disk would be pointless load. Messages, which must be
durable, go through \texttt{postgres\_changes}. Two primitives, two purposes.}

\Q{How do the unread badges stay correct?}
\A{A trigger. \texttt{on\_message\_inserted} fires after every insert and increments
\texttt{host\_unread} or \texttt{guest\_unread} on the parent conversation depending
on \texttt{sender\_role}, and stores a 100-character preview. The client never
increments a counter, so the count is right no matter how the message was inserted.}

\Q{Why does a sent message appear before the network call finishes?}
\A{Optimistic UI. We generate a UUID client-side, append a message with
\texttt{sending: true}, and render images from a local
\texttt{URL.createObjectURL} blob so they appear instantly. Then we upload in
parallel and insert the real row --- \emph{using the same id}. When the realtime
event for our own message comes back, the \texttt{prev.some(m => m.id ===
newMsg.id)} check recognises it and does not duplicate. If the upload fails we
remove the optimistic message and revoke the blob URLs.}

\Q{Why revoke the object URLs?}
\A{\texttt{URL.createObjectURL} pins the file in memory until you revoke it. In a
long chat session that is a genuine memory leak, so we revoke on both the success
and the failure path.}

\Q{How do you compress images?}
\A{Client-side, before upload, with the Canvas API. We draw the image into a canvas
scaled so its longest side is at most 1200 pixels, then call
\texttt{canvas.toBlob(..., "image/jpeg", 0.6)}. A 12~MB photo becomes about
150~KB. It saves the user's bandwidth, our storage bill, and the recipient's load
time.}

\Q{The guest can see the host's phone number. Doesn't RLS say you can only read
your own profile?}
\A{Yes, and that is exactly the problem we had to solve. Rather than weakening the
policy, we wrote a \texttt{SECURITY DEFINER} function
\texttt{shares\_conversation\_with(profile\_id)} that returns true only if an
actual conversation exists between the caller and that profile, and added it as a
second SELECT policy. So contact details unlock precisely when a real conversation
exists --- never before. Public pages do not touch \texttt{profiles} at all; they
read the \texttt{public\_profiles} view.}

\subsection{QR check-in}

\Q{What is inside the QR code?}
\A{A URL: \texttt{https://.../bookings/<uuid>/check-in}. Generated server-side with
the \texttt{qrcode} package as a base64 data URI, so nothing is stored or uploaded.
Because it is a URL, a guest scanning it with their normal camera app lands on
their own check-in page.}

\Q{So if I photograph someone's QR code I can check in as them?}
\A{No. The id is not a credential. To act on it the scanner must call
\texttt{check\_in\_booking()}, which asks
\texttt{can\_operate\_hotel(booking.hotel\_id)} --- true only for that hotel's
manager or assigned staff. A stranger cannot even read the booking, because RLS
blocks the select.}

\Q{What stops a booking being checked in twice?}
\A{\texttt{check\_in\_booking} requires \texttt{status = 'confirmed'}. After the
first check-in the status is \texttt{'checked\_in'}, so the second attempt raises
``Only confirmed bookings can be checked in''. The state machine is the guard.}

\Q{How does the camera scanning work?}
\A{The \texttt{html5-qrcode} library requests the rear camera
(\texttt{facingMode: "environment"}) and decodes 10 frames a second. We pull the
UUID out of the scanned URL with a regex, falling back to the raw text so a bare
booking id also works. The \texttt{useEffect} cleanup calls \texttt{stop()} ---
without that, the camera light stays on after you close the dialog.}

\subsection{Maps, email, and the rest}

\Q{Why not Google Maps?}
\A{It needs an API key that must be shipped to the browser, where anyone can lift
it and spend your money, and it needs a credit card on file. MapLibre with
OpenFreeMap vector tiles is fully open source, needs no key, and has no rate limit.
We kept Leaflet with CartoDB raster tiles as an automatic fallback so the map never
disappears if WebGL is unavailable.}

\Q{What is the difference between vector and raster tiles?}
\A{Raster tiles are pre-rendered PNG images --- zoom in and they blur. Vector tiles
ship geometry and the browser draws it with WebGL, so labels stay crisp at any
zoom and we can restyle the map (that muted ``positron'' palette) entirely
client-side.}

\Q{Why Brevo, and why the HTTP API rather than SMTP?}
\A{Serverless functions are short-lived and frequently blocked from outbound SMTP
ports, whereas an HTTPS \texttt{fetch} always works. Brevo's free tier also does not
require owning a verified domain. And \texttt{sendEmail} returns \texttt{false}
rather than throwing --- a mail outage must never fail a paid booking.}

\Q{Why is the email HTML written with tables and inline styles?}
\A{Because email clients strip \texttt{<style>} blocks and Outlook renders with
Word's layout engine, which has no flexbox or grid. Tables plus inline styles are
the only combination that renders identically in Gmail, Outlook and Apple Mail.}

\Q{What is the difference between a manager and staff?}
\A{A manager owns hotels and must be approved by an admin before they can list
anything. Staff are real login accounts, invited by a manager, scoped to \emph{one}
hotel through a \texttt{hotel\_staff} row carrying a permissions array. Every
database check goes through \texttt{can\_operate\_hotel()} (manager or any staff)
or \texttt{can\_manage\_hotel(hotel, perm)} (manager or staff with that specific
permission).}

\Q{Could a staff member accept an invite meant for someone else?}
\A{No. \texttt{accept\_my\_staff\_invites()} matches pending invites against the
\emph{authenticated user's email}, not just the token in the link. A forwarded link
redeemed by the wrong account matches nothing. Invites also expire after 14 days.}

\Q{How do you stop an admin from silently doing something?}
\A{Every admin power is a \texttt{SECURITY DEFINER} RPC that begins with
\texttt{if get\_my\_role() <> 'admin' then raise exception} and ends by inserting
into \texttt{admin\_audit\_log} --- in the same transaction as the change. An
unlogged admin action is not possible.}

\Q{How do you show prices in different currencies?}
\A{Prices are stored in INR because that is what Razorpay charges. At display time
\file{lib/currency.ts} detects the visitor's currency from their IP (falling back to
the browser locale), fetches live exchange rates, and formats with
\texttt{Intl.NumberFormat}. Both the rates and the detected locale are cached in
\texttt{localStorage} --- six hours and 24 hours respectively. If anything fails we
fall back to showing INR. \textbf{Conversion is presentation only}; the charge is
always the stored INR amount.}

\subsection{The hard questions}

\Q{What is the weakest part of this system?}
\A{Three things, honestly. First, \file{next.config.ts} sets
\texttt{images.unoptimized = true} and allows \texttt{hostname: "**"}, because a
manager may paste any image URL. The comment in the file says we should tighten it
once all images live in Supabase Storage. Second, the \texttt{manager-documents}
storage bucket allows anonymous inserts, because the document must be uploaded
before the account exists --- the schema comment already flags that a production
system should move this behind a server action. Third, the older
\texttt{confirm\_payment} RPC (the manual UPI path, from before Razorpay) is
trust-based: it takes the guest's word for a UPI reference number.}

\Q{What would you build next?}
\A{Automated tests around \texttt{book\_room} --- a concurrency test that fires N
simultaneous bookings at the last room and asserts exactly one succeeds. Then rate
limiting on the auth and booking endpoints. Then moving the manager-document upload
behind a signed server action.}

\Q{What was the hardest bug?}
\A{The RLS infinite recursion between \texttt{hotels} and \texttt{hotel\_staff}.
The error message points at a policy, but the cause is a \emph{cycle} between two
policies, so nothing looks wrong in either one individually. The fix --- single-table
\texttt{SECURITY DEFINER} helpers --- is the entire content of
\file{supabase/hms\_rls\_recursion\_fix.sql}.}

\Q{If you had to explain your project's architecture in one sentence?}
\A{A Next.js app where the security and the business rules live in Postgres --- Row
Level Security decides what you can see, and PL/pgSQL functions decide what you can
do --- so the front end can be careless and the system is still correct.}

\Q{Why should the database, not the API, enforce security?}
\A{Because an API has many entry points and each one is a chance to forget a check.
A table has one set of policies. If I add a new endpoint tomorrow and forget to
check ownership, RLS still returns nothing. Defence in depth, with the deepest
layer doing the real work.}

\newpage
%==============================================================================
\section{Known Limitations and Honest Notes}

Being able to name your own weaknesses is worth more than pretending there are
none. If asked ``what's wrong with it?'', pick from this list.

\begin{enumerate}[leftmargin=*]
  \item \textbf{No automated tests.} The concurrency logic in \texttt{book\_room}
        deserves a test that fires simultaneous bookings at the last room.
  \item \textbf{Images are unoptimised.} \texttt{next.config.ts} disables Next's
        image optimiser and allows any HTTPS host, because managers may paste
        arbitrary image URLs. The mitigation --- an Unsplash-specific URL rewriter
        --- lives in \file{lib/image.ts}, but it only helps for Unsplash.
  \item \textbf{The \texttt{manager-documents} bucket accepts anonymous uploads},
        because the document is uploaded before the account exists. The schema
        comment already flags this: \emph{``For a production system, move this
        upload behind a server action.''}
  \item \textbf{The legacy \texttt{confirm\_payment} RPC is trust-based.} It was
        the manual-UPI path before Razorpay. It is still granted to
        authenticated users and should be revoked.
  \item \textbf{No rate limiting} on auth or booking endpoints.
  \item \textbf{Currency conversion is display-only} and depends on a third-party
        FX API. If it is down, everything falls back to INR.
  \item \textbf{A typo in \file{supabase/booking\_schema.sql} line 23}
        --- \texttt{num\_rooms int not nu ll default 1} --- would make that file fail
        if run on a fresh database. The live database was created before the typo
        was introduced, so it has never been noticed. It should be fixed before
        anyone tries to reproduce the schema from scratch.
\end{enumerate}

\begin{notebox}
That last one is a real bug in the repository. Fixing it before the demo takes ten
seconds, and mentioning that you found it while auditing your own schema is a much
better story than having someone else find it.
\end{notebox}

\newpage
%==============================================================================
\section{Appendix A: Environment Variables}

\begin{longtable}{@{}p{6.2cm}p{2.0cm}p{6.4cm}@{}}
\toprule
\textbf{Variable} & \textbf{Exposed?} & \textbf{Purpose} \\
\midrule
\endhead
\texttt{NEXT\_PUBLIC\_SUPABASE\_URL} & Browser & Which Supabase project. \\
\texttt{NEXT\_PUBLIC\_SUPABASE\_ANON\_KEY} & Browser & Public key. Safe --- RLS
guards the data. \\
\texttt{SUPABASE\_SERVICE\_ROLE\_KEY} & \textbf{Server only} &
\textbf{Bypasses RLS.} Webhook only. \\
\texttt{NEXT\_PUBLIC\_SITE\_URL} & Browser & Auth redirect target. \\
\texttt{RAZORPAY\_KEY\_ID} & Server & Razorpay API id. \\
\texttt{RAZORPAY\_KEY\_SECRET} & \textbf{Server only} & Signs and verifies
payments. \\
\texttt{NEXT\_PUBLIC\_RAZORPAY\_KEY\_ID} & Browser & Publishable key for the
checkout modal. \\
\texttt{RAZORPAY\_WEBHOOK\_SECRET} & \textbf{Server only} & Verifies the webhook
body HMAC. \\
\texttt{BREVO\_API\_KEY} & \textbf{Server only} & Sends transactional email. \\
\texttt{EMAIL\_FROM}, \texttt{EMAIL\_FROM\_NAME} & Server & Sender identity. \\
\bottomrule
\end{longtable}

\begin{warnbox}
\textbf{The rule.} In Next.js, \emph{any} variable prefixed
\texttt{NEXT\_PUBLIC\_} is inlined into the JavaScript bundle at build time and is
therefore readable by anyone. Everything without that prefix is server-only. Notice
that \texttt{RAZORPAY\_KEY\_ID} and \texttt{NEXT\_PUBLIC\_RAZORPAY\_KEY\_ID} hold
the same value --- one for server use, one deliberately marked as public. The
\emph{secret} never gets the prefix.
\end{warnbox}

%==============================================================================
\section{Appendix B: Glossary}
\label{sec:glossary}

\begin{longtable}{@{}p{4.0cm}p{10.8cm}@{}}
\toprule
\textbf{Term} & \textbf{Meaning} \\
\midrule
\endhead
\textbf{App Router} & Next.js's routing system where folders are URLs. \\
\textbf{Atomic} & All-or-nothing. Either every step happens or none does. \\
\textbf{CDN} & A network of servers worldwide that cache your pages close to users. \\
\textbf{Compare-and-swap} & Change a value only if it still holds the expected
value. Used for our once-only email. \\
\textbf{HMAC} & A fingerprint of a message made with a shared secret. Proves who
sent it and that it was not altered. \\
\textbf{Idempotent} & Doing it twice has the same effect as doing it once. \\
\textbf{ISR} & Incremental Static Regeneration --- a page is pre-rendered and
refreshed on a timer. \\
\textbf{JWT} & JSON Web Token. A signed blob proving who you are. \\
\textbf{Optimistic UI} & Showing the result before the server confirms it, and
rolling back if it fails. \\
\textbf{Paise} & 1/100 of a rupee. Razorpay works in integer paise to avoid
floating-point money bugs. \\
\textbf{PostgREST} & The service that turns Postgres tables into a REST API. \\
\textbf{Race condition} & Two operations interleaving in a way that corrupts data. \\
\textbf{Row lock} & \texttt{SELECT ... FOR UPDATE}. Makes other transactions wait
for that row. \\
\textbf{RLS} & Row Level Security. Per-row access rules enforced by Postgres. \\
\textbf{RPC} & Remote Procedure Call. Here: calling a Postgres function from the
app. \\
\textbf{RSC} & React Server Component. Runs on the server, ships no JavaScript. \\
\textbf{SECURITY DEFINER} & A Postgres function that runs with its creator's
privileges, bypassing RLS internally. \\
\textbf{Server Action} & A server function a Client Component can call directly. \\
\textbf{SWR} & Stale-While-Revalidate. Show cached data now, refresh in background. \\
\textbf{Timing attack} & Learning a secret by measuring how long a comparison takes. \\
\textbf{Trigram index} & An index over every 3-character substring, making
\texttt{ILIKE '\%x\%'} fast. \\
\textbf{WebSocket} & A connection that stays open so the server can push data. \\
\bottomrule
\end{longtable}

\vspace{1.2cm}
\begin{center}
\rule{0.5\textwidth}{0.6pt}\\[8pt]
{\large\color{primary}\textbf{BookNest}}\\[3pt]
{\small Built by Team Neuron}\\[10pt]
\emph{Good luck tomorrow. You know this system --- now you can prove it.}
\end{center}

\end{document}
